\documentclass[11pt,oneside,openany]{book}
\usepackage[T1]{fontenc}
\usepackage{lmodern,microtype}
\usepackage[a4paper,margin=28mm]{geometry}
\usepackage{amsmath,amssymb,amsthm,mathtools}
\usepackage[hidelinks,bookmarksopen,bookmarksnumbered]{hyperref}
\hypersetup{pdftitle={Foundations of Elementary Number Theory},
  pdfauthor={MathGenome},pdfsubject={Unified notes, P2-03}}
\allowdisplaybreaks[1]
\setlength{\emergencystretch}{2em}

\theoremstyle{plain}
\newtheorem{theorem}{Theorem}[chapter]
\newtheorem{lemma}[theorem]{Lemma}
\newtheorem{proposition}[theorem]{Proposition}
\newtheorem{corollary}[theorem]{Corollary}
\theoremstyle{definition}
\newtheorem{definition}[theorem]{Definition}
\newtheorem{example}[theorem]{Example}
\theoremstyle{remark}
\newtheorem{remark}[theorem]{Remark}

\begin{document}
\frontmatter
\hypersetup{pageanchor=false}
\begin{titlepage}
\thispagestyle{empty}
\vspace*{27mm}
{\large MathGenome \enspace / \enspace P2-03\par}
\vspace{8mm}
\hrule height 0.8pt
\vspace{11mm}
{\Huge\bfseries Foundations of\\[2mm] Elementary Number Theory\par}
\vspace{12mm}
{\large Unified notes\par}
\vspace{6mm}
{\normalsize Integer arithmetic, congruences, and reciprocity\\
Arithmetic functions and elementary prime estimates\\
Diophantine equations, continued fractions, and approximation\par}
\vfill
{\large 8 November 2026\par}
\end{titlepage}
\hypersetup{pageanchor=true}
\tableofcontents

\chapter*{Scope and conventions}
\addcontentsline{toc}{chapter}{Scope and conventions}
These notes develop the arithmetic of integers and finite residue rings,
followed by arithmetic functions, sums of squares, and rational
approximation. The final chapter uses continued fractions to determine
the positive solutions of both signs of Pell's equation. The treatment
draws on Hardy and Wright for the broad organization, Apostol for
arithmetic functions and congruences, and Niven, Zuckerman, and Montgomery
for the Diophantine and approximation chapters
\cite{hardy-wright,apostol,nzm}.

Moduli and rational denominators are positive unless explicitly stated
otherwise. A prime is a positive prime integer, and gcds are nonnegative.
Prime valuations of integers are used on nonzero arguments unless a
zero convention is given. Arithmetic functions have positive-integer
arguments, and a divisor sum runs over positive divisors. The two
normalizations for rational approximation are distinguished throughout:
the error of a fraction and the error of its integer numerator after
multiplying by its denominator.

The standard prerequisites are elementary real analysis (limits, series,
differentiation and integration), elementary group and field theory
(including splitting fields), and finite-dimensional linear algebra.
The three-square proof uses ordinary Euclidean volume. Its auxiliary
prime comes from Dirichlet's theorem on primes in reduced residue
classes, proved in Chapter~\ref{ch:nt-diophantine}. That chapter proves
the required character identities and nonvanishing at \(1\), the
lattice argument, and the conversion of its output into three
integer squares.

\mainmatter
\chapter{Divisibility and integer congruences}
\label{chap:integers}

This chapter fixes the arithmetic conventions used throughout the notes.
We write $\mathbb N=\{0,1,2,\ldots\}$, $\mathbb N_{>0}$ for the positive
integers, and $a\mid b$ if $b=ac$ for an integer $c$. In particular,
$0\mid b$ holds precisely when $b=0$. Gcds and lcms are nonnegative.
All congruence moduli are positive, and $a\equiv b\pmod m$ means
$m\mid(a-b)$. See \cite[Chs.~I, V, VIII]{hardy-wright},
\cite[Chs.~1 and 5]{apostol}, and \cite[Chs.~2 and 4]{shoup}.

\section{Euclidean division and linear combinations}

\begin{lemma}[Division with remainder]\label{nt:division}
For $a\in\mathbb Z$ and $m>0$ there are unique integers $q,r$ such that
$a=qm+r$ and $0\le r<m$.
\end{lemma}
\begin{proof}
Choose the least nonnegative member $r$ of $\{a-km:k\in\mathbb Z\}$.
This set has a nonnegative member by taking $k$ sufficiently negative,
and $r\ge m$ would make $r-m$ a smaller one. Two possible remainders
differ by a multiple of $m$ strictly between $-m$ and $m$, so they
agree and their quotients agree.
\end{proof}

\begin{theorem}[Euclidean algorithm, B\'ezout, and Euclid's lemma]
\label{nt:A01}
Given $a,b\in\mathbb Z$, the Euclidean algorithm applied to $|a|,|b|$
returns their nonnegative gcd $d$. The set $a\mathbb Z+b\mathbb Z$
equals $d\mathbb Z$, and every common divisor of $a,b$ divides $d$.
For $(a,b)=(0,0)$ this reads $d=0$. If a prime $p$ divides $ab$, then
$p$ divides $a$ or $b$. The least common multiple of $a,b$ is $0$
when $ab=0$ and $|ab|/d$ otherwise. In all cases the common multiples
are exactly the multiples of the lcm.
\end{theorem}
\begin{proof}
Start with $r_0=|a|$ and $r_1=|b|$. Whenever $r_i>0$, divide
$r_{i-1}=q_i r_i+r_{i+1}$ with $0\le r_{i+1}<r_i$. At each
step
\[
 r_{i-1}\mathbb Z+r_i\mathbb Z=r_i\mathbb Z+r_{i+1}\mathbb Z.
\]
One inclusion uses the remainder formula, and the other uses the
division formula. Strict descent forces a final pair $(d,0)$. If
$b=0$ the initial pair already has this form. The initial ideal
equals $a\mathbb Z+b\mathbb Z$, so the displayed invariance proves the
claim about linear combinations. In particular $d$ divides both
integers and is a linear combination of them. This also gives the
extended algorithm: retain the coefficients of $r_0,r_1$ as
combinations of $a,b$ and update the coefficient vector of
$r_{i+1}$ by subtracting $q_i$ times that of $r_i$ from that of
$r_{i-1}$. The coefficient vector of the final $d$ is a B\'ezout
pair.

If $p\nmid a$, the only positive common divisor of $p,a$ is $1$, so
$up+va=1$ for some integers $u,v$. Multiplication by $b$ shows that
$p\mid ab$ implies $p\mid b$. More generally, if $\gcd(x,y)=1$
and $x\mid yz$, the identity $ux+vy=1$ multiplied by $z$ proves
$x\mid z$.

For $ab\ne0$, write $|a|=dA$, $|b|=dB$. Dividing the B\'ezout
identity by $d$ shows $\gcd(A,B)=1$. The positive number $dAB$
is a common multiple. If $M$ is any common multiple, write
$M=dAu$ and use $B\mid Au$ and coprime cancellation to deduce
$B\mid u$. This gives $dAB\mid M$. When $ab=0$, the only
common multiple is $0$.
\end{proof}

\section{Prime factorization and valuations}

A prime is an integer $p>1$ whose only positive divisors are $1,p$.
For a nonzero integer $n$, its sign is $+1$ or $-1$ according to
whether $n>0$ or $n<0$. The units of $\mathbb Z$ are exactly
$1,-1$, since $xy=1$ forces $|x|=|y|=1$.

\begin{theorem}[Unique factorization and prime valuations]\label{nt:A02}
Every positive integer is a product of primes, unique up to order,
with the empty product for $1$. Each nonzero integer is its sign
times such a product. For a prime $p$ and $n\ne0$, define $v_p(n)$
as the exponent of $p$ in this factorization. For nonzero $a,b$,
\[
 v_p(ab)=v_p(a)+v_p(b),\qquad
 a\mid b\ \Longleftrightarrow\ v_p(a)\le v_p(b)
                       \text{ for every prime }p.
\]
For nonzero $a,b$, the $p$-valuations of their gcd and lcm are the
minimum and maximum of $v_p(a),v_p(b)$, respectively.
\end{theorem}
\begin{proof}
Use strong induction for existence. A composite $n>1$ has a
factorization $n=uv$ with $1<u,v<n$, so its factors have prime
factorizations. For uniqueness, a prime in one product divides
the other product. Repeated Euclid's lemma gives a factor on that
side which it divides, and primality forces the factors to agree.
Cancel that factor and repeat. This procedure cannot leave a
nonempty prime product equal to $1$. Taking absolute values and
restoring the sign proves the integer version.

For nonzero inputs, multiplying products adds the exponents. If
$a\mid b$, write $b=ac$ and apply this formula, since $c\ne0$.
Conversely, use the differences of the exponents and the sign of
$b$ divided by the sign of $a$ to construct $c\in\mathbb Z$ with
$b=ac$. The numbers with the componentwise minimum and maximum
exponents, and positive sign, therefore satisfy the defining gcd
and lcm divisibility properties.
\end{proof}

In particular, $v_p(1)=v_p(-1)=0$ and $p^j\mid n$ exactly when
$j\le v_p(n)$ for $n\ne0$. One may extend $v_p$ by $v_p(0)=+\infty$
in $\mathbb N\cup\{+\infty\}$. Then the product formula uses
$+\infty+k=+\infty$, and divisibility still has the stated
comparison interpretation. Indeed, every $a$ divides $0$, and
$0$ divides only $0$. With this convention the gcd and lcm
formulas also hold when either input is $0$, using
$\min(+\infty,+\infty)=\max(+\infty,+\infty)=+\infty$.
Any use of a finite valuation explicitly assumes the input is
nonzero. A frequently used consequence is that if $v_p(U)=j$
is finite and $p^{j+1}\mid V$, then $v_p(U+V)=j$. Dividing by
$p^j$ gives a nonzero residue plus zero modulo $p$.

\section{Linear equations and linear congruences}

\begin{theorem}[Linear Diophantine equations and congruences]\label{nt:A03}
Let $a,b,c\in\mathbb Z$ and $d=\gcd(a,b)$. The equation $ax+by=c$
has an integer solution exactly when $d\mid c$. If $d>0$ and
$(x_0,y_0)$ is one solution, all solutions are
\[
 (x,y)=(x_0+(b/d)t,\;y_0-(a/d)t),\qquad t\in\mathbb Z.
\]
For $m>0$ the congruence $ax\equiv b\pmod m$ is solvable exactly
when $d'=\gcd(a,m)$ divides $b$. If solvable, it has exactly $d'$
solutions as classes modulo $m$.
\end{theorem}
\begin{proof}
By Theorem~\ref{nt:A01}, the set of possible values of $ax+by$
is $d\mathbb Z$. When $d=0$ both coefficients vanish, so every
pair is a solution for $c=0$ and there are no solutions otherwise.
For $d>0$, subtract a fixed solution and divide by $d$. The
differences $X=x-x_0$ and $Y=y-y_0$ obey
$(a/d)X+(b/d)Y=0$ with coprime coefficients. If $b\ne0$, coprime
cancellation gives $X=(b/d)t$, and substitution gives $Y=-(a/d)t$.
If $b=0$, then $a/d=\pm1$, so $X=0$ and every $Y$ is again of the
specified form. Substitution proves that every pair in the formula
is a solution.

The congruence is equivalent to $ax+my=b$ for some $y\in\mathbb Z$.
When $d'\mid b$, the preceding formula shows that $x$ ranges over
$x_0+(m/d')\mathbb Z$. Two parameters $t,t'$ give the same class
modulo $m$ exactly when $d'\mid(t-t')$. Thus the $d'$ parameters
$0,\ldots,d'-1$ yield all distinct classes, including the case
$m=1$.
\end{proof}

\section{Factorials and binomial coefficients}

\begin{theorem}[Legendre's factorial formula]\label{nt:A04}
If $p$ is prime and $n\in\mathbb N$, then
\[
 v_p(n!)=\sum_{j\ge1}\left\lfloor\frac{n}{p^j}\right\rfloor .
\]
For $0\le k\le n$,
\[
 v_p\binom nk
   =\sum_{j\ge1}\left(
       \left\lfloor\frac n{p^j}\right\rfloor
       -\left\lfloor\frac k{p^j}\right\rfloor
       -\left\lfloor\frac{n-k}{p^j}\right\rfloor
     \right).
\]
Every summand in the latter sum is $0$ or $1$, and all summands
with $p^j>n$ vanish.
\end{theorem}
\begin{proof}
Use additivity of $v_p$ and count the pairs $(i,j)$ with
$1\le i\le n$ and $1\le j\le v_p(i)$. For fixed $j$, the number
of choices for $i$ is the number of positive multiples of $p^j$
at most $n$, which is $\lfloor n/p^j\rfloor$. When $n=0$ both
counts are empty and $0!=1$. The binomial coefficient counts the
$k$-element subsets of an $n$-element set. Counting permutations
by their first $k$ elements gives
$n!=\binom nk k!(n-k)!$, so additivity gives the second formula.
For each $j$, write $k$ and $n-k$ by division with remainder
modulo $p^j$. The difference of the floors is the floor of the
sum of these two remainders divided by $p^j$. The sum belongs to
$[0,2p^j)$, as required.
\end{proof}
This elementary counting formula is used for prime-counting
bounds. See \cite[\S8.1]{nzm}.

\section{The Chinese remainder theorem}

\begin{theorem}[Finite Chinese remainder theorem]\label{nt:A05}
Let $m_1,\ldots,m_s>0$ and $a_1,\ldots,a_s\in\mathbb Z$. There is
an integer $x$ satisfying all $x\equiv a_i\pmod{m_i}$ if and only
if
\[
 a_i\equiv a_j\pmod{\gcd(m_i,m_j)}\qquad(1\le i,j\le s).
\]
If a solution exists, its class is unique modulo
$L=\operatorname{lcm}(m_1,\ldots,m_s)$. If the moduli are pairwise
coprime and $M=\prod_i m_i$, the reduction map is a ring
isomorphism
\[
 \mathbb Z/M\mathbb Z \longrightarrow
     \prod_{i=1}^s \mathbb Z/m_i\mathbb Z.
\]
For $m_i>1$, set $M_i=M/m_i$, choose $u_i$ with
$M_i u_i\equiv1\pmod{m_i}$, and set $e_i=M_i u_i$. The inverse
sends $(a_i)_i$ to the class of $\sum_{m_i>1}a_i e_i$ modulo $M$.
The empty lcm and product are both $1$.
\end{theorem}
\begin{proof}
Necessity follows from the divisibility of both $x-a_i$ and
$x-a_j$ by their moduli. For two moduli $m,n$, write $x=a+mt$.
The remaining condition is $mt\equiv b-a\pmod n$, so
Theorem~\ref{nt:A03} proves sufficiency. The difference of any
two solutions is divisible by both $m,n$, hence by their lcm,
and adding an lcm multiple preserves both congruences.

For more moduli, suppose a solution $y$ is known for the first
$s-1$ conditions. Write $L'$ for the lcm of the first $s-1$
moduli. For each prime power $p^k$ dividing $\gcd(L',m_s)$, some
previous $m_i$ is divisible by $p^k$, by the componentwise maximum
formula for $v_p(L')$. The hypotheses give $p^k\mid(a_i-a_s)$,
and $p^k\mid(y-a_i)$ by the choice of $y$. Thus $p^k\mid(y-a_s)$.
As this holds for each prime power in $\gcd(L',m_s)$, their product
divides the difference. Apply the two-modulus case to
$y\bmod L'$ and $a_s\bmod m_s$. The difference and translation
argument with the full lcm proves uniqueness. A one-modulus
system needs no induction, and for the empty system every
integer is a solution, giving one class modulo $1$.

For pairwise coprime moduli, the lcm equals the product by the
valuation formulas. Reduction is a well-defined ring homomorphism
as each $m_i$ divides $M$. Its bijectivity follows from existence
and uniqueness. The cofactor $M_i$ is coprime to $m_i$, so
B\'ezout gives $u_i$. Each $e_i$ is $1$ in its own component
and $0$ in every other, proving the inverse formula. A modulus
$1$ contributes a one-element factor and imposes no condition.
The same formula includes the empty product as the one-element
ring.
\end{proof}
Compare \cite[Thm.~5.26]{apostol}. The finite coprime isomorphism,
existence, uniqueness, and explicit inverse already have
counterparts in \cite[D07, Ch.~9]{mathgenome}.

\section{Two elementary infinitude results}

\begin{theorem}[Infinitely many primes]\label{nt:A06}
There are infinitely many primes. There are also infinitely many
primes congruent to $3$ modulo $4$.
\end{theorem}
\begin{proof}
For a proposed finite list of all primes, take their product $P$,
where $P=1$ if the list is empty. The integer $P+1>1$ has a prime
factor by Theorem~\ref{nt:A02}. A prime factor of $P$ cannot
divide $P+1$, so the list is incomplete.

For a proposed finite list of primes congruent to $3\bmod4$,
take their product $P$ and factor $4P-1>1$. This number is odd
and is $3\bmod4$. Every odd prime is $1$ or $3$ modulo $4$, and
a product made entirely from primes congruent to $1$ would itself
be $1$. Thus at least one prime factor is $3\bmod4$. It cannot
be on the list, since the listed primes divide $4P$.
\end{proof}
The progression result here uses only factorization. A later
application to three squares invokes the general prime-existence
theorem for reduced arithmetic progressions, proved in
Theorem~\ref{nt:dirichlet-primes}.
For the elementary prime arguments, see \cite[Chs.~I and II]{hardy-wright}
and \cite[\S\S1.4, 1.6, 7.2]{apostol}.
\chapter{Residue groups and quadratic reciprocity}
\label{chap:congruences}

For $m\ge1$, write $U_m=(\mathbb Z/m\mathbb Z)^\times$ for the
multiplicative group of units. A primitive root modulo $m$ is a
representative of a generator of $U_m$. This convention includes
$m=1$, whose sole class is its multiplicative identity. The
proofs in this chapter follow the arithmetic conventions of
Chapter~\ref{chap:integers}. Main references are
\cite[Chs.~5, 9, 10]{apostol},
\cite[Chs.~2 and 3]{nzm}, and \cite[Ch.~12]{shoup}.

\section{Euler, Fermat, and Wilson}

For an element $g$ of a finite group, its order is the least positive
$d$ such that $g^d=1$. A positive power with this property exists
because two among $1,g,g^2,\ldots$ agree. Division of any integer
exponent by $d$ shows that $g^j=1$ exactly when $d\mid j$.
The powers $1,\ldots,g^{d-1}$ comprise a subgroup of size $d$.
Its left cosets partition the group into sets of this size, so
$d$ divides the group order.

\begin{theorem}[Totient and classical congruences]\label{nt:B01}
Set $\phi(m)=|U_m|$, equivalently the number of $0\le a<m$
with $\gcd(a,m)=1$. Then $\phi(1)=1$, $\phi(p^k)=p^{k-1}(p-1)$
for every prime $p$ and $k\ge1$, and
\[
 \phi(m)=m\prod_{p\mid m}\left(1-\frac1p\right),\qquad
 \phi(mn)=\phi(m)\phi(n)\quad\text{if }\gcd(m,n)=1.
\]
If $\gcd(a,m)=1$ then $a^{\phi(m)}\equiv1\pmod m$. In particular,
for a prime $p$ and $p\nmid a$ one has $a^{p-1}\equiv1\pmod p$.
For every integer $a$ and prime $p$ one has $a^p\equiv a\pmod p$.
Finally, $m>1$ is prime exactly when $(m-1)!\equiv-1\pmod m$.
\end{theorem}
\begin{proof}
B\'ezout gives an inverse whenever $\gcd(a,m)=1$. Conversely,
$ab\equiv1\pmod m$ gives $ab+mc=1$ for an integer $c$, so the
gcd is $1$. This applies also to the one class modulo $1$.
For $p^k$, among the $p^k$ representatives there are $p^{k-1}$
multiples of $p$, and the others are units. The isomorphism in
Theorem~\ref{nt:A05} restricts to units because multiplication
and inversion in a product are componentwise. This proves the
totient formulas by unique factorization.

The order of a unit divides $|U_m|$, by the coset argument.
This gives Euler's congruence. For $m=p$ multiply the resulting
unit congruence by $a$ to obtain the all-residue formula when
$p\nmid a$. If $p\mid a$, both sides vanish modulo $p$.

For odd prime $p$, the only self-inverse units are $1$ and $-1$.
Indeed $p\mid(x-1)(x+1)$ forces one factor to vanish modulo
$p$, by Euclid's lemma. All other units pair with a different
inverse, contributing $1$ per pair, while $1$ and $-1$ contribute
$-1$. For $p=2$, the factorial is $1\equiv-1\pmod2$.
If $m>1$ is composite, choose a divisor $d$ with $1<d<m$.
It divides $(m-1)!$, so the factorial cannot be congruent to
$-1$ even modulo $d$. Congruence modulo $m$ would imply
congruence modulo $d$.
\end{proof}
See \cite[\S\S2.3--2.5 and Thms.~5.17, 5.19, 5.24]{apostol}.

\section{Finite fields and cyclic multiplicative groups}

A nonzero polynomial of degree $d$ over a field has at most
$d$ distinct roots. Division by $X-\alpha$ shows that each
root $\alpha$ gives such a factor, and induction on degree
gives the bound. We also use the following algebraic construction.

\begin{lemma}[Splitting fields]\label{nt:splitting-field}
For a nonconstant polynomial over a field $K$, there is a field
generated over $K$ by all of its roots in which the polynomial
splits, unique up to $K$-isomorphism.
\end{lemma}
\begin{proof}
If the polynomial has a nonlinear irreducible factor $h$, adjoining
the class of $X$ in $K[X]/(h)$ gives a field containing a root of
$h$. Irreducible factors exist by induction on degree, and the
quotient is a field because polynomial division and the Euclidean
algorithm give a B\'ezout identity between $h$ and every
nonzero residue class. The number of linear factors, counted
with multiplicity, increases when adjoining a root of a
nonlinear factor. Repeating in the enlarged fields eventually
makes the polynomial split. The subfield generated over $K$
by the resulting roots is the desired splitting field.

For uniqueness, let $L,L'$ be two splitting fields. Starting with
the identity on $K$, extend an embedding into $L'$ one root at
a time. If a root $\alpha$ has irreducible polynomial $h$ over
the currently embedded subfield, then $h$ divides the original
polynomial over that subfield. This follows by division by $h$,
since a nonzero remainder would have smaller degree and vanish
at $\alpha$. The image of $h$ has a root in $L'$ because it divides
a split polynomial. The quotient construction extends the embedding
by sending $\alpha$ to that root. Once all roots have been
adjoined, the image contains a full linear factorization of the
polynomial, hence all of its roots in $L'$. They generate $L'$.
\end{proof}

\begin{theorem}[Finite fields and cyclicity]\label{nt:B02}
Every finite field has prime characteristic $p$ and cardinality
$p^f$ for an integer $f\ge1$. For each prime $p$ and $f\ge1$
there is a field of this cardinality, unique up to an isomorphism
fixing the prime field. Every finite subgroup of the multiplicative
group of a field is cyclic. In particular, a field of cardinality
$q$ has a cyclic multiplicative group of order $q-1$ with exactly
$\phi(q-1)$ generators.
\end{theorem}
\begin{proof}
In a finite field some positive multiple of $1$ is zero, by
finiteness of the additive group. Its least positive value $p$
must be prime. If $p=ab$ with $1<a,b<p$, then
$(a\cdot1)(b\cdot1)=0$ with both factors nonzero. Thus the field
contains the image of $\mathbb F_p=\mathbb Z/p\mathbb Z$. As a
vector space over it the field has a finite basis, obtained by
adjoining elements outside the span until it spans the finite
set. A basis of size $f$ gives $p^f$ coordinate vectors.

Conversely put $q=p^f$ and take a splitting field of $X^q-X$
over $\mathbb F_p$. The formal derivative is $-1$, so no root
can be repeated, since a repeated linear factor would make the
derivative vanish there. In characteristic $p$ the binomial
coefficients $\binom pj$ for $0<j<p$ are divisible by $p$
(by Theorem~\ref{nt:A04}), so the map $x\mapsto x^p$ is additive
and multiplicative. Iteration proves the same for $x\mapsto x^q$.
Its fixed points are therefore closed under addition,
subtraction, and multiplication. They contain $0,1$ and are
closed under inverses of nonzero elements, because
$(x^{-1})^q=(x^q)^{-1}$. This root set is a field with $q$
elements. Every field of order $q$ is a splitting field of the
same polynomial, since every nonzero element satisfies
$x^{q-1}=1$ by the finite group argument and all $q$ elements
are roots. Lemma~\ref{nt:splitting-field} proves uniqueness.

Let $H$ be a finite subgroup of the multiplicative group of
a field, and let $e$ be the lcm of its element orders. For
each prime power $\ell^a$ occurring exactly in $e$, some element
of $H$ has order $\ell^a t$ with $\ell\nmid t$. Its $t$-th
power has order $\ell^a$, by the criterion for exponents
annihilating an element. Two commuting elements of coprime
orders $r,s$ have product of order $rs$. In fact a power that
kills the product makes the corresponding powers of the two
elements inverse, so each belongs to the intersection of
their cyclic subgroups. An element of this intersection has
order dividing both $r$ and $s$ and is the identity. Therefore
both $r$ and $s$ divide the exponent. Taking the product of the
chosen prime-power-order elements yields an element of order
$e$, with the identity for $e=1$. All of $H$ consists of roots
of $X^e-1$, so $|H|\le e$ by the root bound, whereas the element
of order $e$ has $e$ distinct powers in $H$. Thus $H$ is cyclic.

Finally, if $g$ has order $d$, then $g^j$ has order
$d/\gcd(d,j)$, since $d\mid jt$ exactly when
$d/\gcd(d,j)\mid t$. The generators are the $\phi(d)$ classes of
exponents coprime to $d$. Apply this to $d=q-1$, including
$q=2$, when $d=1$.
\end{proof}
See \cite[Ch.~2]{lidl} and \cite[\S10.4]{apostol}.
The cardinality, existence, uniqueness, and unit-group
statements have separate counterparts in \cite[D07, Ch.~13]{mathgenome}.

\section{Primitive roots for positive moduli}

\begin{theorem}[Lifting a primitive root at an odd prime]\label{nt:B03}
Let $p$ be an odd prime. If $u\in\mathbb Z$, $u\equiv1\pmod p$,
and $u\ne1$, then
\[
 v_p(u^p-1)=v_p(u-1)+1.
\]
There is an integer $g$ that generates $U_{p^k}$ for every
$k\ge1$.
\end{theorem}
\begin{proof}
Write $u=1+p^j w$ with $j\ge1$ and $p\nmid w$. The linear term
of $(1+p^jw)^p-1$ is $p^{j+1}w$. For $2\le i<p$, the
$i$-th term is divisible by $p^{1+ij}$, hence by $p^{j+2}$.
The final term is divisible by $p^{jp}$, which is also at
least $p^{j+2}$ since $p\ge3$. Thus after division by $p^{j+1}$
the residue is $w\pmod p$, giving the exact valuation.

Choose a representative $g_0$ of a generator of
$\mathbb F_p^\times$ with $1\le g_0<p$. Its order $p-1$ is at
least $2$, so $g_0\ge2$. Both $g_0^{p-1}$ and
$(g_0+p)^{p-1}$ are $1\bmod p$, and their difference modulo
$p^2$ is
\[
  p(p-1)g_0^{p-2}.
\]
This is nonzero modulo $p^2$. At least one of $g_0,g_0+p$
therefore satisfies $v_p(g^{p-1}-1)=1$. The preceding valuation
formula, iterated from $u=g^{p-1}>1$, gives
\[
 v_p\!\left(g^{(p-1)p^r}-1\right)=1+r,\qquad r\ge0.
\]
The order of $g$ modulo $p^k$ is a multiple of $p-1$ by
reduction modulo $p$, and divides $(p-1)p^{k-1}$ by
Theorem~\ref{nt:B01}. For $k\ge2$, the displayed valuation
at $r=k-2$ excludes every proper divisor of $(p-1)p^{k-1}$
that is a multiple of $p-1$. Thus $g$ has order $(p-1)p^{k-1}$
and generates the unit group. The case $k=1$ holds by choice.
\end{proof}

\begin{theorem}[Classification and count of primitive roots]\label{nt:B04}
For $m\ge1$, $U_m$ is cyclic exactly for
\[
 m=1,\ 2,\ 4,\ p^k,\ \text{or }2p^k
 \qquad(p\text{ an odd prime},\ k\ge1).
\]
For each cyclic $U_m$, the number of primitive-root classes is
$\phi(\phi(m))$. For $k\ge3$, the units modulo $2^k$ have the
unique form $(-1)^a5^b$ with $a\in\{0,1\}$ and
$0\le b<2^{k-2}$.
\end{theorem}
\begin{proof}
For each $r\ge0$,
\[
 v_2\!\left(5^{2^r}-1\right)=r+2.
\]
For $r=0$ this is $v_2(4)=2$. In the induction step,
factor the next difference as
$(5^{2^r}-1)(5^{2^r}+1)$. The second factor is $2\bmod4$,
so the valuation increases by one. The order of $5$ modulo
$2^k$ divides the group order $2^{k-1}$ and hence is a power
of two. The first $r$ with $5^{2^r}\equiv1\pmod{2^k}$ is
$k-2$, so the order is $2^{k-2}$. Its powers lie in the
$1\bmod4$ classes, and those classes number $2^{k-2}$, so
the powers exhaust them. Each unit is $1$ or $3\bmod4$.
A unit in the latter class is the negative of one in the former.
This proves existence and uniqueness of the displayed form.
It also produces distinct elements of order two, namely $-1$
and $5^{2^{k-3}}$, for every $k\ge3$. A cyclic group has
at most one element of order two: in a cyclic group of
order $d$, the only possible nonidentity solution of
$2j\equiv0\pmod d$ is $j=d/2$ when $d$ is even.

The groups $U_1,U_2$ are trivial and $U_4$ is cyclic of order $2$.
Theorem~\ref{nt:B03} treats $p^k$, and $U_{2p^k}$ is
isomorphic to $U_{p^k}$ by the Chinese remainder theorem since
$U_2$ is trivial. For the remaining cases write
$m=2^j\prod_i p_i^{k_i}$ with distinct odd primes $p_i$.
If $j\ge3$, the two order-two elements of $U_{2^j}$ give
two in the product. If there are at least two odd-prime
factors, choose the element $-1$ in each of two distinct
factors and the identity elsewhere to give two such elements.
If $j=2$ and there is an odd-prime factor, use $-1$ in $U_4$
and $-1$ in an odd-prime factor in the same way. These cases
exhaust the remaining $m$. Finally the count follows from
the generator count in Theorem~\ref{nt:B02} applied to a
cyclic group of order $\phi(m)$.
\end{proof}
For the lift and classification, see
\cite[\S\S10.1--10.9 and Thm.~10.11]{apostol}. The classification
has a fixed-version counterpart in \cite[UnitsCyclic]{mathlib}.

\section{Simple-root lifting}

\begin{theorem}[Polynomial roots modulo prime powers]\label{nt:B05}
Let $p$ be prime, $F\in\mathbb Z[X]$, $j\ge1$, and
$a\in\mathbb Z$ with $F(a)\equiv0\pmod{p^j}$ and
$p\nmid F'(a)$. Among the $p$ classes
$a+p^j t\pmod{p^{j+1}}$ with $t\bmod p$, exactly one is a
root of $F$ modulo $p^{j+1}$. It has the same nonzero derivative
residue modulo $p$ as $a$. Consequently a simple root modulo $p$
extends uniquely at every higher prime power.
\end{theorem}
\begin{proof}
For a monomial, the binomial expansion at $a+p^jt$ has constant
term, linear term, and all terms of order at least two divisible
by $p^{2j}$. Summing the monomials gives
\[
 F(a+p^jt)\equiv F(a)+p^jtF'(a)\pmod{p^{j+1}}
\]
because $2j\ge j+1$. Dividing the required congruence by $p^j$
leads to the linear congruence
\[
 tF'(a)\equiv-F(a)/p^j\pmod p,
\]
which has one solution modulo $p$ since $F'(a)$ is a unit
there. Every lift of $a\bmod p^j$ has the indicated form,
and $F'(a+p^jt)\equiv F'(a)\pmod p$. Induction proves the
full lifting assertion.
\end{proof}
This is the finite-modulus simple-root result. For a root with
zero derivative modulo $p$, a separate analysis is needed.
See \cite[Thms.~5.28 and 5.30]{apostol}.

\section{Legendre symbols and Gauss's lemma}

For an odd prime $p$, the Legendre symbol $\left(\frac ap\right)$
is $0$ if $p\mid a$, $1$ if $a$ represents a nonzero square
modulo $p$, and $-1$ otherwise.

\begin{theorem}[Euler's criterion and Gauss's lemma]\label{nt:B06}
Put $h=(p-1)/2$. For all integers $a,b$,
\[
 a^h\equiv\left(\frac ap\right)\pmod p,\qquad
 \left(\frac{ab}p\right)=\left(\frac ap\right)\left(\frac bp\right).
\]
If $p\nmid a$, let $N$ count the least positive residues of
$a,2a,\ldots,ha$ that are greater than $p/2$. Then
$\left(\frac ap\right)=(-1)^N$.
\end{theorem}
\begin{proof}
The zero case of Euler's criterion is immediate since $h\ge1$.
For a unit $a$, write it as a power $g^j$ of a generator of
$\mathbb F_p^\times$. The squares are exactly the even powers,
since the group order $2h$ is even. The element $g^h$ has order
two, so it is $-1$ by the factorization $(X-1)(X+1)$ over the
field. Thus $a^h=(-1)^j$, giving the criterion. Multiplicativity
follows because the congruences respect products and the three
integers $-1,0,1$ are distinct modulo $p$.

For Gauss's lemma, replace each residue of $ia$ above $p/2$
by its negative signed representative in $[-h,-1]$, and keep the
others in $[1,h]$. The absolute values are all distinct. If the
values for $i,j$ agree then $ia\equiv\pm ja\pmod p$, giving
$i\equiv\pm j\pmod p$. Equality with a plus sign forces $i=j$;
the minus sign is impossible since $2\le i+j\le p-1$. Thus
the absolute values are a permutation of $1,\ldots,h$.
Multiplying the signed representatives gives
$a^h h!\equiv(-1)^N h!\pmod p$. The factorial is invertible
modulo $p$, and Euler's criterion identifies the resulting sign.
\end{proof}
See \cite[\S\S9.1--9.4]{apostol} and \cite[\S\S3.1--3.2]{nzm}.

\section{Quadratic reciprocity and its supplements}

\begin{theorem}[Quadratic reciprocity]\label{nt:B07}
For distinct odd primes $p,q$,
\[
 \left(\frac qp\right)\left(\frac pq\right)
   =(-1)^{((p-1)/2)((q-1)/2)}.
\]
For every odd prime $p$,
\[
 \left(\frac{-1}p\right)=(-1)^{(p-1)/2},\qquad
 \left(\frac 2p\right)=(-1)^{(p^2-1)/8}.
\]
\end{theorem}
\begin{proof}
The first supplement follows from Euler's criterion with $a=-1$.
For the second, the least residues of $2i$ for
$1\le i\le h=(p-1)/2$ are exactly $2,4,\ldots,2h$.
The number greater than $p/2$ is $h-\lfloor p/4\rfloor$.
It is even for $p\equiv1,7\pmod8$ and odd for
$p\equiv3,5\pmod8$. This follows by writing $p=8t+r$ with
$r=1,3,5,7$. The same parities hold for $(p^2-1)/8$, since
it is congruent to $(r^2-1)/8$ modulo $2$. Gauss's lemma gives
the result.

For the main law, first let $a$ be any positive odd integer
coprime to $p$ and write
$ia=p\lfloor ia/p\rfloor+r_i$ with $1\le r_i\le p-1$ for
$1\le i\le h$. The absolute values of the signed representatives
are a permutation of $1,\ldots,h$. Changing a residue
$r_i>p/2$ to its signed representative changes its absolute-value
parity, because $r_i$ and $p-r_i$ have opposite parity.
If $N$ is the number of such changes, then
$\sum_i r_i\equiv\sum_i i+N\pmod2$. Reducing the division
identities modulo $2$ and using $a,p$ odd gives
\[
 \sum_{i=1}^h\left\lfloor\frac{ia}{p}\right\rfloor\equiv N\pmod2.
\]
This gives the floor-sum form of Gauss's lemma.

Let $k=(q-1)/2$. In the rectangle of pairs
$1\le i\le h$, $1\le j\le k$, there is no pair with $q i=p j$,
since it would force $p\mid i$. For fixed $i$, the number
of pairs with $p j<q i$ is $\lfloor qi/p\rfloor$. Indeed
$qi/p$ is not an integer and is less than $q/2$, so its
floor is at most $k$. By the same argument for fixed $j$,
the complementary count is $\lfloor pj/q\rfloor$.
Thus the sum of the two floor sums is $hk$, and the two
floor-sum signs multiply to $(-1)^{hk}$.
\end{proof}
See \cite[\S\S9.3--9.5]{apostol}, \cite[\S\S3.1--3.2]{nzm},
and \cite[Ch.~12]{shoup}.

\section{The Jacobi symbol and its computation}

For positive odd $m=\prod_p p^{e_p}$, define
\[
 \left(\frac am\right)_{\!J}
   =\prod_{p:e_p>0}\left(\frac ap\right)^{e_p},
\]
where the factors on the right are Legendre symbols. The empty
product for $m=1$ is $1$. When $m$ is an odd prime, the two
symbols agree.

\begin{theorem}[Jacobi rules]\label{nt:B08}
The Jacobi symbol depends only on the class of the numerator
modulo its denominator, is multiplicative in each argument, and
is zero exactly when the arguments have gcd greater than $1$.
Its supplement formulas are
\[
 \left(\frac{-1}m\right)_{\!J}=(-1)^{(m-1)/2},\qquad
 \left(\frac 2m\right)_{\!J}=(-1)^{(m^2-1)/8}.
\]
For positive odd $a,m$,
\[
 \left(\frac am\right)_{\!J}
   =(-1)^{((a-1)/2)((m-1)/2)}\left(\frac ma\right)_{\!J}.
\]
These rules compute the symbol by Euclidean reductions. A value
$+1$ alone does not imply a square residue for composite $m$.
\end{theorem}
\begin{proof}
The assertions about the residue class, multiplication, and
zero follow from the prime factors and Theorem~\ref{nt:B06}.
For positive odd $u,v$ we have
\begin{align*}
 \frac{uv-1}{2}&\equiv\frac{u-1}{2}+\frac{v-1}{2}\pmod2,\\
 \frac{(uv)^2-1}{8}&\equiv\frac{u^2-1}{8}+\frac{v^2-1}{8}\pmod2.
\end{align*}
In the first identity the omitted term is $(u-1)(v-1)/2$,
which is even. In the second the omitted term is
$(u^2-1)(v^2-1)/8$, a multiple of $8$ because odd squares
are $1\bmod8$. These identities combine the prime-factor
supplements into the stated formulas.

If $a,m$ are not coprime, both symbols in the reciprocity
formula vanish. Otherwise all prime-factor pairs are distinct.
Apply Theorem~\ref{nt:B07} to each pair, with multiplicities.
The total sign exponent is the product of
$\sum_{p\mid a} v_p(a)(p-1)/2$ and
$\sum_{q\mid m} v_q(m)(q-1)/2$. The first parity identity
reduces these sums to $(a-1)/2$ and $(m-1)/2$ modulo $2$.
This also covers a denominator or numerator equal to $1$.

For the computation, first return $1$ if $m=1$. Otherwise
reduce the numerator to $0\le r<m$. If $r=0$, return $0$.
For $r>0$, extract the power of two $r=2^v c$ with $c$ odd,
and remove its symbol by the second supplement. If $c=1$,
the remaining symbol is $1$. For $c>1$, reciprocity gives
the known sign and the symbol with numerator $m$ and
denominator $c$. Repeat with that smaller denominator.
This terminates and uses only the proved equalities, including
when a zero result occurs. For example
$\left(\frac{2}{15}\right)_{\!J}
 =\left(\frac23\right)\left(\frac25\right)=1$,
but $2$ is not a square modulo $3$ and hence not modulo $15$.
\end{proof}
See \cite[\S9.7]{apostol} and \cite[\S3.3]{nzm}.

\section{All roots of a quadratic congruence}

\begin{theorem}[Quadratic congruences over arbitrary positive moduli]
\label{nt:B09}
Let $a\in\mathbb Z$, $m>0$, and write $m=\prod_p p^{k_p}$ over
the positive prime exponents. The number of roots of
$x^2\equiv a\pmod m$ is the product of the root counts below
for the prime powers $p^{k_p}$. The empty product is $1$.
For $p^k$, $k\ge1$, the count is given as follows.
\begin{enumerate}
 \item If $a\equiv0\pmod{p^k}$, the roots are the classes
       divisible by $p^{\lceil k/2\rceil}$, so there are
       $p^{\lfloor k/2\rfloor}$ roots.
 \item Otherwise $s=v_p(a)<k$. If $s$ is odd there are no
       roots. If $s=2t$, set $j=k-2t$ and $a'=a/p^{2t}$.
       The count is $p^t$ times the number of roots of
       $y^2\equiv a'\pmod{p^j}$, where $p\nmid a'$.
       For odd $p$, this latter count is $2$ if
       $\left(\frac{a'}p\right)=1$ and $0$ otherwise.
       For $p=2$, the unit counts are $1$ for $j=1$,
       $2$ for $j=2$ when $a'\equiv1\pmod4$, and $4$ for
       $j\ge3$ when $a'\equiv1\pmod8$. In the last two
       ranges the count is zero if the indicated condition
       fails.
\end{enumerate}
\end{theorem}
\begin{proof}
The Chinese remainder bijection identifies the equation with
the independent equations in all prime-power components,
so it multiplies the counts. If $a$ is the zero class modulo
$p^k$, the condition is $2v_p(x)\ge k$ for nonzero $x$.
The same divisibility criterion includes $x=0$. The number
of classes divisible by $p^{\lceil k/2\rceil}$ is the quotient
$p^k/p^{\lceil k/2\rceil}=p^{\lfloor k/2\rfloor}$.

If $a$ is not the zero class, changing its representative by
a multiple of $p^k$ preserves its valuation $s<k$. The
same fact applied to $x^2=a+p^k z$ shows that a root must
have $2v_p(x)=s$, excluding odd $s$. When $s=2t$, write
$x=p^t y$. The classes $x\bmod p^k$ divisible by $p^t$ are
in bijection with $y\bmod p^{k-t}$. After dividing the
congruence by $p^{2t}$ the required equation is
$y^2\equiv a'\pmod{p^{k-2t}}$, with $y$ automatically a unit.
Every class of $y$ modulo $p^{k-2t}$ has exactly $p^t$ lifts
modulo $p^{k-t}$, and all these lifts give distinct $x$.

For odd $p$, a square unit modulo $p$ has two distinct roots
$\pm y$, by the degree-two root bound, and a nonsquare has
none. Each root of $Y^2-a'$ is simple since $2y\not\equiv0\pmod p$,
and Theorem~\ref{nt:B05} lifts it uniquely. This gives the
odd-prime count. Modulo $2$ there is one odd class, and modulo
$4$ both odd classes square to $1$. If $j\ge3$, any odd
square is $1\bmod8$: writing an odd integer as $2r+1$ gives
its square as $1+4r(r+1)$, and $r(r+1)$ is even.
In the decomposition from Theorem~\ref{nt:B04}, squaring kills
the sign and doubles the exponent of $5$ modulo $2^{j-2}$.
Its kernel has size $4$, since the exponent has two choices
with $2b\equiv0\pmod{2^{j-2}}$ and the sign is arbitrary.
Its image has size $2^{j-1}/4=2^{j-3}$, because every
nonempty fibre is a coset of the kernel. This equals the
number of classes congruent to $1\bmod8$, so the image
is precisely those classes, each with four roots.
\end{proof}
The count separates singular roots from simple roots and includes
the degenerate modulus $1$. The prime-power ingredients can be
compared with \cite[\S\S5.9, 9.6, 10.11]{apostol}.
\chapter[Arithmetic functions and primes]{Arithmetic Functions and Elementary Prime Estimates}
\label{ch:arithmetic-functions}

Throughout this chapter \(\mathbb N_{>0}\) denotes the positive integers.
In a sum with a real upper limit \(x\), an index \(n\le x\) always means a
positive integer \(n\) satisfying the displayed inequality. All logarithms
are natural logarithms. Statements in \(O\)-notation or \(\asymp\)-notation
refer to \(x\to+\infty\). We write \(F\asymp G\) for two eventual inequalities
\(cG\le F\le CG\) with constants \(c,C>0\).

\section{Dirichlet convolution and local factors}

\begin{theorem}[Convolution and its units]\label{nt:C01}
Let \(R\) be a commutative ring with identity. The functions
\(f\colon\mathbb N_{>0}\to R\) form a commutative ring under pointwise
addition and the product
\[
 (f*g)(n)=\sum_{d\mid n}f(d)g(n/d).
\]
The identity \(\delta\) takes value \(1\) at \(1\) and \(0\) elsewhere.
A function is a unit for this product exactly when \(f(1)\) is a unit in \(R\).
\end{theorem}

\begin{proof}
Write the product as a sum over ordered pairs \((a,b)\) with \(ab=n\).
Interchanging \(a\) and \(b\) gives commutativity, and expanding either
association of a triple product gives
\(\sum_{abc=n}f(a)g(b)h(c)\), with each ordered triple occurring once.
Distributivity follows from distributivity in \(R\). The term with index
\(1\) in the \(\delta\) factor gives the identity law.

If \(f*g=\delta\), evaluation at \(1\) gives \(f(1)g(1)=1\).
Conversely, for \(f(1)\) a unit, define \(g(1)=f(1)^{-1}\) and for \(n>1\) set
\[
 g(n)=-f(1)^{-1}\sum_{\substack{d\mid n\\d>1}} f(d)g(n/d).
\]
Each \(n/d\) on the right is strictly less than \(n\), so the recursion
defines \(g\) on all positive integers. The equations it imposes are
precisely the coefficients of \(f*g=\delta\).
\end{proof}

\begin{theorem}[Multiplicativity and Euler products]\label{nt:C02}
A function \(f\colon\mathbb N_{>0}\to R\) is called multiplicative if
\(f(1)=1\) and \(f(ab)=f(a)f(b)\) whenever \(\gcd(a,b)=1\). Such a function
is determined by its prime-power values through
\[
 f(n)=\prod_{p\mid n} f\bigl(p^{v_p(n)}\bigr).
\]
Any assignment of values at positive prime powers extends uniquely to a
multiplicative function. Multiplicative functions are closed under
convolution and inversion in the convolution ring.

If \(f\) is complex-valued and multiplicative and a real \(s\) satisfies
\(\sum_{n\ge1}|f(n)|n^{-s}<\infty\), then
\[
 \sum_{n\ge1}\frac{f(n)}{n^s}
   =\prod_p\left(\sum_{j\ge0}\frac{f(p^j)}{p^{js}}\right).
\]
The product is the limit of the products over \(p\le y\) as \(y\to\infty\).
In particular, for \(s>1\),
\(\zeta(s)=\sum_{n\ge1} n^{-s}=\prod_p(1-p^{-s})^{-1}\).
\end{theorem}

\begin{proof}
Unique prime factorization and induction on the number of distinct prime
divisors give the product formula. For the converse, coprime factors
have disjoint sets of prime divisors, so the formula is multiplicative.
The empty product gives the prescribed value at \(1\).

If \(\gcd(a,b)=1\), each divisor of \(ab\) is uniquely \(d_1d_2\), where
\(d_1\mid a\), \(d_2\mid b\). Both \(d_1,d_2\) and \(a/d_1,b/d_2\) are
coprime pairs. Multiplicativity of \(f,g\) therefore factors the divisor
sum for \((f*g)(ab)\) as \((f*g)(a)(f*g)(b)\).
For the inverse, set \(g(1)=1\) and assign the prime-power values by
\[
 g(p^j)=-\sum_{i=1}^j f(p^i)g(p^{j-i})\qquad(j\ge1).
\]
Extend these values multiplicatively. By construction \(f*g\) vanishes
at every positive prime power and is \(1\) at \(1\). Since \(f*g\) is
multiplicative, it equals \(\delta\).

For the analytic assertion, each local series converges absolutely as
its terms form a subseries of the given absolutely convergent series.
A finite product of these series expands absolutely into the sum over
integers all of whose prime factors lie among the finitely many selected
primes. Each coefficient follows from multiplicativity and each integer
in that sum has a unique exponent tuple. Selecting \(p\le y\), every
omitted integer exceeds \(y\). The difference from the full series is
bounded by \(\sum_{n>y}|f(n)|n^{-s}\), which tends to zero. For the
constant function \(1\) and \(s>1\), convergence follows by comparison
with \(\int_1^\infty t^{-s}\,dt\), and the local geometric series is
\((1-p^{-s})^{-1}\).
\end{proof}

The finite convolution and the analytic passage to prime products use
different hypotheses. For the source treatments of both, see
\cite[Sections 2.6--2.11, 2.15, and 11.2--11.5]{apostol} and
\cite[Chapters XVI--XVII]{hardy-wright}.

\section{Classical arithmetic functions and M\"obius inversion}

\begin{proposition}[Classical functions]\label{nt:C03}
Write \(\mathbf{1}(n)=1\), \(\mathrm{id}_k(n)=n^k\), and
\(\log(n)=\log n\) for arithmetic functions. The M\"obius function \(\mu\)
vanishes on integers with a squared prime factor and has value
\((-1)^r\) on a product of \(r\) distinct primes. The Euler totient
\(\phi(n)\) counts the integers from \(1\) to \(n\) coprime to \(n\).
The divisor count \(\tau(n)\) and the divisor power sum \(\sigma_k(n)\)
satisfy, for \(k\ge0\),
\[
 \tau=\mathbf{1}*\mathbf{1},\qquad
 \sigma_k=\mathbf{1}*\mathrm{id}_k .
\]
These four functions are multiplicative. For primes \(p\) and \(j\ge1\),
\[
 \mu(p^j)=
 \begin{cases}-1& j=1,\\0&j\ge2,\end{cases}
 \quad
 \phi(p^j)=p^j-p^{j-1},\quad
 \tau(p^j)=j+1,\quad
 \sigma_k(p^j)=\sum_{i=0}^j p^{ik}.
\]
In particular \(\sigma_0=\tau\). The von Mangoldt function
\(\Lambda(n)=\log p\) when \(n=p^j\) with \(j\ge1\) and \(0\) otherwise
satisfies \(\mathbf{1}*\Lambda=\log\).
\end{proposition}

\begin{proof}
The empty product gives \(\mu(1)=1\). The definition of \(\mu\) on
prime-power factors and Theorem~\ref{nt:C02} give multiplicativity.
The integers from \(1\) to \(n\) represent each class modulo \(n\) once.
Thus the residue-class description of \(\phi\) in Theorem~\ref{nt:B01}
gives multiplicativity via the Chinese remainder theorem and the
prime-power formula by deleting the \(p^{j-1}\) multiples of \(p\)
among \(p^j\) classes. Both \(\mathbf{1}\) and \(\mathrm{id}_k\) are multiplicative,
and the two divisor sums and their prime-power formulas follow by
listing the divisors. Finally, if \(n=\prod_p p^{e_p}\), the only
divisors that contribute to \(\sum_{d\mid n}\Lambda(d)\) are
\(p,p^2,\ldots,p^{e_p}\). The resulting sum is
\(\sum_p e_p\log p=\log n\), which is \(0\) at \(n=1\).
\end{proof}

\begin{theorem}[M\"obius inversion]\label{nt:C04}
For functions \(f,g\colon\mathbb N_{>0}\to R\), with integer-valued
coefficients sent to \(R\) by the unital map \(\mathbb Z\to R\), one has
\[
 g(n)=\sum_{d\mid n}f(d)\ \text{for all }n
 \quad\Longleftrightarrow\quad
 f(n)=\sum_{d\mid n}\mu(d)g(n/d)\ \text{for all }n .
\]
In the integer-valued convolution ring,
\(\mu*\mathbf{1}=\delta\), \(\sum_{d\mid n}\phi(d)=n\), and
\(\phi=\mu*\mathrm{id}_1\). In the real-valued ring
\(\Lambda=\mu*\log\).
\end{theorem}

\begin{proof}
The multiplicative function \(\mu*\mathbf{1}\) takes value
\(\mu(1)+\mu(p)+\cdots+\mu(p^j)=0\) at \(p^j\) for \(j\ge1\).
Its value at \(1\) is \(1\), so it is \(\delta\). Mapping the integer
coefficients into \(R\) gives the same identity there. Multiplication by
\(\mu\) sends \(g=\mathbf{1}*f\) to \(f=\mu*g\), and multiplication by
\(\mathbf{1}\) reverses this step.

For each \(d\mid n\), the integers \(k\) with \(1\le k\le n\) and
\(\gcd(k,n)=d\) are the \(k=dt\) with \(1\le t\le n/d\) and
\(\gcd(t,n/d)=1\). Thus \(n=\sum_{d\mid n}\phi(n/d)
=\sum_{d\mid n}\phi(d)\). Applying inversion gives the totient formula.
The Mangoldt formula follows by applying inversion to
\(\mathbf{1}*\Lambda=\log\) over \(\mathbb R\).
\end{proof}

See \cite[Sections 2.2--2.8 and 2.13]{apostol} and
\cite[Sections 4.2--4.3]{nzm} for these functions and identities.

\section{Exact summation identities}

\begin{theorem}[Dirichlet hyperbola identity]\label{nt:C05}
For \(f,g\colon\mathbb N_{>0}\to R\), put \(F(t)=\sum_{n\le t}f(n)\)
and \(G(t)=\sum_{n\le t}g(n)\). Let \(U,V\) be nonnegative integers and
\(x\) a real number with \(UV\le x<(U+1)(V+1)\). Then
\[
 \sum_{n\le x}(f*g)(n)=
 \sum_{u\le U}f(u)G(x/u)+\sum_{v\le V}g(v)F(x/v)-F(U)G(V).
\]
In particular one can take \(U=V=\lfloor\sqrt{x}\rfloor\).
\end{theorem}

\begin{proof}
The sum on the left is the sum of \(f(u)g(v)\) over positive integer
pairs with \(uv\le x\). Every such pair has \(u\le U\) or \(v\le V\).
Indeed a pair with \(u\ge U+1\) and \(v\ge V+1\) has \(uv>x\).
The two sums on the right count these two subsets. Their overlap is
exactly the rectangle \(u\le U,\ v\le V\), since its pairs satisfy
\(uv\le UV\le x\). Subtracting the full rectangle once yields the
identity. If \(U=0\) or \(V=0\) the relevant rectangle and sum are empty,
and the same reasoning applies. The symmetric choice of cut satisfies
both required inequalities.
\end{proof}

\begin{theorem}[Partial summation]\label{nt:C06}
Let \(A(t)=\sum_{n\le t}a_n\). For \(1\le m\le N\), the discrete
identity in a commutative ring is
\[
 \sum_{n=m}^{N} a_n w_n
 =A(N)w_N-A(m-1)w_m+\sum_{n=m}^{N-1} A(n)(w_n-w_{n+1}).
\]
For real or complex coefficients, \(a\le b\) real, and a weight
\(w\) of class \(C^1\) on \([a,b]\), the Abel identity is
\[
 \sum_{a<n\le b} a_nw(n)
 =A(b)w(b)-A(a)w(a)-\int_a^b A(t)w'(t)\,dt .
\]
When \(a=b\) this is interpreted with empty sum and zero integral.
\end{theorem}

\begin{proof}
For the discrete version, insert \(a_n=A(n)-A(n-1)\) and shift the
index of the second sum. The two boundary terms and the paired
interior terms are the ones displayed.

For \(a<b\) and \(t\in[a,b]\), the finite decomposition
\(A(t)=A(a)+\sum_{a<n\le b}a_n\mathbf{1}_{[n,b]}(t)\)
gives by the fundamental theorem of calculus
\[
 \int_a^b A(t)w'(t)\,dt
 =A(a)(w(b)-w(a))
   +\sum_{a<n\le b}a_n(w(b)-w(n)).
\]
Here \(\mathbf{1}_{[n,b]}\) is the indicator of the interval. An integer
\(a\) contributes through \(A(a)\), and an integer \(b\) gives the
zero-length integral at \(n=b\). Since
\(A(b)=A(a)+\sum_{a<n\le b}a_n\), rearrangement gives the formula with
the stated endpoints.
\end{proof}

See \cite[Theorems 3.17 and 4.2]{apostol} and
\cite[Chapters 2--3]{gkp}.

\section{Three average orders}

\begin{lemma}[Harmonic numbers and the Basel sum]
For \(H_m=\sum_{j=1}^m1/j\), the limit
\(\gamma=\lim_{m\to\infty}(H_m-\log m)\) exists, and
\(H_m=\log m+\gamma+O(1/m)\). Also
\(\sum_{n\ge1}n^{-2}=\pi^2/6\).
\end{lemma}

\begin{proof}
The difference \((H_m-\log m)-(H_{m+1}-\log(m+1))\) is
\(\int_m^{m+1} dt/t-1/(m+1)\). It is nonnegative and bounded above by
\(1/m-1/(m+1)\). The sequence is bounded below by \(0\), by the integral
comparison \(H_m\ge\log(m+1)\), so it converges to \(\gamma\). Summing
the bounds on the differences from \(m\) to infinity bounds the
difference from \(\gamma\) by \(1/m\).

For the second identity let \(0\le r<1\), \(0\le t\le\pi\). The geometric
series and taking imaginary parts give
\[
 \frac{\partial}{\partial r}
    \sum_{n\ge1}\frac{r^n\sin(nt)}{n}
 =\sum_{n\ge1}r^{n-1}\sin(nt)
 =\frac{\sin t}{1-2r\cos t+r^2}.
\]
Differentiation is justified by uniform convergence on
\(0\le r\le R<1\). The derivative on the right is also the derivative
of \(\arctan(r\sin t/(1-r\cos t))\), and both antiderivatives vanish
at \(r=0\). Integrating the resulting identity in \(t\) gives
\[
 \sum_{n\ge1}\frac{r^n(1-(-1)^n)}{n^2}
 =\int_0^\pi \arctan\!\left(\frac{r\sin t}{1-r\cos t}\right)dt .
\]
The integration of the series uses uniform convergence at each fixed
\(r<1\). As \(r\) increases to \(1\), the integrand tends to
\((\pi-t)/2\) for \(0<t\le\pi\). On \([\varepsilon,\pi]\) the
convergence is uniform because the denominators are bounded below by
\(\min(1,1-\cos\varepsilon)>0\). On \([0,\varepsilon]\) both the
integrand and its limiting expression are bounded by \(\pi/2\), which
allows the integral to converge on the full interval. On the left,
the summable majorant \(2/n^2\) bounds the tails uniformly in \(r\).
The limit is therefore
\[
 2\sum_{n\text{ odd}}\frac1{n^2}=\frac{\pi^2}{4}.
\]
The even terms comprise one quarter of the full reciprocal-square
sum, so the left side is three halves of that sum.
\end{proof}

\begin{theorem}[Divisors, totients, and squarefree integers]\label{nt:C07}
One has
\begin{align*}
 \sum_{n\le x}\tau(n)&=x\log x+(2\gamma-1)x+O(\sqrt{x}),\\
 \sum_{n\le x}\phi(n)&=\frac{3}{\pi^2}x^2+O(x\log x),\\
 \sum_{n\le x}\mu(n)^2&=\frac{6}{\pi^2}x+O(\sqrt{x}).
\end{align*}
\end{theorem}

\begin{proof}
For \(x\ge1\) put \(m=\lfloor\sqrt{x}\rfloor\). Theorem~\ref{nt:C05}
with \(f=g=\mathbf{1}\) gives
\[
 \sum_{n\le x}\tau(n)=2\sum_{d\le m}\lfloor x/d\rfloor-m^2
                  =2xH_m-m^2+O(m).
\]
Here \(m^2=x+O(\sqrt{x})\). Also
\(0\le\log\sqrt{x}-\log m\le 1/m\) and \(m\ge\sqrt{x}/2\).
The harmonic estimate therefore gives the first formula.

For the other two formulas, absolute convergence of both
\(\sum\mu(n)n^{-2}\) and \(\sum n^{-2}\) allows the double sum
of their product to be grouped by the product of its indices. By
Theorem~\ref{nt:C04} this gives
\[
 \left(\sum_{d\ge1}\frac{\mu(d)}{d^2}\right)
 \left(\sum_{m\ge1}\frac{1}{m^2}\right)
 =\sum_{n\ge1}\frac{(\mu*\mathbf{1})(n)}{n^2}=1 .
\]
Thus \(\sum_{d\ge1}\mu(d)/d^2=6/\pi^2\).
The totient convolution yields
\[
 \sum_{n\le x}\phi(n)
 =\sum_{d\le x}\mu(d)\sum_{m\le x/d}m
 =\frac{x^2}{2}\sum_{d\le x}\frac{\mu(d)}{d^2}
       +O\!\left(x\sum_{d\le x}\frac1d\right).
\]
For \(y=x/d\ge1\), the error in replacing
\(\lfloor y\rfloor(\lfloor y\rfloor+1)/2\) by \(y^2/2\) is \(O(y)\),
uniformly in \(y\). The harmonic estimate bounds the total error by
\(O(x\log x)\). Extending the main sum to infinity costs \(O(x)\)
because \(\sum_{d>x}d^{-2}=O(1/x)\).

Finally, \(\mu(n)^2=\sum_{d^2\mid n}\mu(d)\). Both sides factor over
prime powers, since a square divisor is specified by independently
choosing its exponent at each prime. At \(p^e\), for \(e=0,1\) the
sum is \(1\), and for \(e\ge2\) it is \(1-1=0\). Thus
\[
 \sum_{n\le x}\mu(n)^2
 =\sum_{d\le\sqrt{x}}\mu(d)\lfloor x/d^2\rfloor
 =x\sum_{d\le\sqrt{x}}\frac{\mu(d)}{d^2}+O(\sqrt{x}).
\]
Extending the main sum costs \(O(\sqrt{x})\), again by the
reciprocal-square tail estimate, proving the third formula.
\end{proof}

The divisor and totient estimates are the classical formulas in
\cite[Theorems 3.3 and 3.7]{apostol}. The squarefree count is treated
in \cite[Theorem 8.25]{nzm}.

\section{Chebyshev estimates and reciprocal primes}

For \(x\ge1\), define
\begin{align*}
 P(x)&=\prod_{p\le x}p,
 &\theta(x)&=\sum_{p\le x}\log p,\\
 \psi(x)&=\sum_{p^k\le x,\ k\ge1}\log p,
 &\pi(x)&=\#\{p\le x:p\text{ prime}\}.
\end{align*}
In this section and the next one, an index \(p\) in a product or sum
over primes always ranges over positive primes.

\begin{lemma}[A bound for the prime product]\label{nt:prime-product-bound}
For every real \(x\ge1\), \(P(x)\le4^x\).
\end{lemma}

\begin{proof}
First take integer \(x=n\) and use strong induction. For \(n=1,2\)
the claim follows from \(P(1)=1\), \(P(2)=2\). For even \(n>2\),
\(n\) is composite, so \(P(n)=P(n-1)\le4^{n-1}\).
For odd \(n=2m+1\ge3\), every prime \(m+1<p\le 2m+1\) divides
\(\binom{2m+1}{m}\). Such a prime divides the numerator factorial
and neither denominator factorial. Distinct primes are coprime,
so their product divides the binomial coefficient. The two equal
coefficients \(\binom{2m+1}{m}\) and \(\binom{2m+1}{m+1}\) have sum
at most \(2^{2m+1}\), hence either one is at most \(4^m\).
Consequently \(P(2m+1)\le P(m+1)4^m\le4^{2m+1}\).
For real \(x\), apply the integer bound to \(\lfloor x\rfloor\).
\end{proof}

\begin{theorem}[Chebyshev's estimates]\label{nt:C08}
One has \(\psi(x)-\theta(x)=O(\sqrt{x})\),
\(\theta(x)\asymp x\), \(\psi(x)\asymp x\), and
\(\pi(x)\asymp x/\log x\).
\end{theorem}

\begin{proof}
The prime product bound gives \(\theta(x)\le x\log4\). Separating
the powers in \(\psi\) gives
\begin{align*}
 0\le\psi(x)-\theta(x)
 &=\theta(\sqrt{x})+
   \sum_{3\le k\le \lfloor\log x/\log 2\rfloor}\theta(x^{1/k})\\
 &=O(\sqrt{x}+x^{1/3}\log x)=O(\sqrt{x}).
\end{align*}
The last comparison follows from the boundedness of
\((\log x)x^{-1/6}\) for \(x\ge2\).

For an integer \(m\ge1\), the factorial valuation formula
\ref{nt:A04} shows that each power \(p^j\le2m\) contributes
\(\lfloor 2m/p^j\rfloor-2\lfloor m/p^j\rfloor\), either \(0\)
or \(1\), to \(v_p\!\left(\binom{2m}{m}\right)\). Therefore
\(\log\binom{2m}{m}\le\psi(2m)\). The central coefficient is the
largest of the \(2m+1\) binomial coefficients, since the ratio of
successive coefficients is \((2m-j)/(j+1)\) for indices \(j,j+1\).
Its logarithm is at least \(m\log4-\log(2m+1)\), since the
coefficients sum to \(4^m\). Taking \(m=\lfloor x/2\rfloor\) and
using monotonicity proves an eventual positive linear lower bound
for \(\psi(x)\). Subtracting \(\psi(x)-\theta(x)=O(\sqrt{x})\)
gives the same lower bound, with a smaller constant, for \(\theta(x)\).
The upper bounds follow from the prime product and power estimates.

For \(x>1\), \(\theta(x)\le \pi(x)\log x\), which gives the lower
bound for \(\pi(x)\). The primes exceeding \(\sqrt{x}\) each contribute
at least \(\tfrac12\log x\) to \(\theta(x)\), and there are at most
\(\sqrt{x}\) smaller primes. Thus
\(\pi(x)\le\sqrt{x}+2\theta(x)/\log x=O(x/\log x)\).
\end{proof}

\begin{theorem}[Two elementary Mertens estimates]\label{nt:C09}
There is a real constant \(B\) such that
\[
 \sum_{p\le x}\frac{\log p}{p}=\log x+O(1),
 \qquad
 \sum_{p\le x}\frac1p=\log\log x+B+O(1/\log x).
\]
In particular, the sum of the reciprocals of the primes diverges.
\end{theorem}

\begin{proof}
Sum \(\log n=\sum_{d\mid n}\Lambda(d)\) over \(n\le x\) to get
\[
 \log(\lfloor x\rfloor!)=\sum_{d\le x}\Lambda(d)\lfloor x/d\rfloor.
\]
Comparing the sum of \(\log n\) for \(n\le N=\lfloor x\rfloor\) to
\(\int_1^N\log t\,dt\) gives an error of size at most \(\log N\).
Thus the left side is \(x\log x+O(x)\). Replacing the floors on
the right costs at most
\(\sum_{d\le x}\Lambda(d)=\psi(x)=O(x)\). Dividing by \(x\) gives
\(\sum_{d\le x}\Lambda(d)/d=\log x+O(1)\).
The terms with prime powers of exponent at least two contribute a
bounded amount, since
\[
 \sum_p\sum_{k\ge2}\frac{\log p}{p^k}
 =\sum_p\frac{\log p}{p(p-1)}
 \le 2\sum_{n\ge2}\frac{\log n}{n^2}<\infty.
\]
The final convergence follows by integral comparison.
This proves the first estimate.

For \(t\ge2\) write \(M(t)=\sum_{p\le t}(\log p)/p=\log t+E(t)\).
Here \(E\) is bounded on \([2,\infty)\). Apply Theorem~\ref{nt:C06}
with weight \(w(t)=1/\log t\) on \([2,x]\) and add the contribution
\(1/2=M(2)w(2)\) from the lower endpoint. Then
\[
 \sum_{p\le x}\frac1p
 =\frac{M(x)}{\log x}+\int_2^x\frac{M(t)}{t(\log t)^2}\,dt.
\]
The \(\log t\) terms give \(1+\log\log x-\log\log2\). The integral
of \(E(t)/(t(\log t)^2)\) converges absolutely on \([2,\infty)\).
Its tail from \(x\) is \(O(1/\log x)\), since
\(\int_x^\infty dt/(t(\log t)^2)=1/\log x\). The term
\(E(x)/\log x\) has the same bound. Taking
\(B=1-\log\log2+\int_2^\infty E(t)\,dt/(t(\log t)^2)\)
yields the result.
\end{proof}

See \cite[Section 8.1, Theorems 8.3--8.5, Corollary 8.6 and
Theorem 8.8]{nzm}, and
\cite[Sections 3.11, 4.2 and 4.8]{apostol}.

\section{Bertrand's postulate}

\begin{theorem}[Bertrand]\label{nt:C10}
For every real \(x>1\), there is a prime \(p\) with \(x<p<2x\).
\end{theorem}

\begin{proof}
First take an integer \(n\ge1024\) and suppose that there is no prime
strictly between \(n\) and \(2n\). The endpoint \(2n\) is composite.
For \(p\le\sqrt{2n}\), the valuation formula used in
Theorem~\ref{nt:C08} gives \(p^{v_p(\binom{2n}{n})}\le2n\), because the
valuation is at most the number of powers \(p^j\le2n\). There are at
most \(\sqrt{2n}\) such primes. For \(\sqrt{2n}<p\le 2n/3\), the
valuation is at most one and the product of the contributions is at
most \(P(2n/3)\le4^{2n/3}\). For \(2n/3<p\le n\), one has \(p^2>2n\),
\(\lfloor 2n/p\rfloor=2\), and \(\lfloor n/p\rfloor=1\), so the
valuation vanishes. These intervals and the assumed absence of
primes up to \(2n\) above \(n\) give
\[
 \frac{4^n}{2n+1}\le\binom{2n}{n}
       \le (2n)^{\sqrt{2n}}\,4^{2n/3}.
\]
The lower bound uses the largest-coefficient argument of
Theorem~\ref{nt:C08}. Taking logarithms and dividing by \(n\) implies
\[
 \frac{\log4}{3}\le
 \frac{\log(2n+1)}{n}+\sqrt{\frac2n}\log(2n).
\]
Both terms on the right decrease for real \(n\ge1024\). The derivative
of the first has the sign of
\(2n/(2n+1)-\log(2n+1)<0\), by integration of \(1/t\) from \(1\) to
\(2n+1\). The derivative of the second has the sign of
\(1-\tfrac12\log(2n)<0\). At \(n=1024\) the right side is less than
\[
 \frac{(12+46\cdot 11)\log2}{1024}<\frac{\log4}{3},
\]
using \(2049<2^{12}\), \(\sqrt{2048}<46\), and \(2048=2^{11}\).
This contradiction proves the claim for \(n\ge1024\).

For \(2\le n<1024\), use the increasing chain
\[
 2,\,3,\,5,\,7,\,13,\,23,\,43,\,83,\,163,\,317,\,631,\,1259 .
\]
Each entry after the first is less than twice the previous entry.
For completeness the rows below list the remainders upon dividing an
entry by the primes at most its square root in increasing order.
A composite number has a prime divisor at most its square root, so
the nonzero remainders certify primality.
\[
\begin{array}{c|l@{\qquad}c|l}
 2&\text{none}&3&\text{none}\\
 5&1&7&1\\
 13&1,1&23&1,2\\
 43&1,1,3&83&1,2,3,6\\
 163&1,1,3,2,9&317&1,2,2,2,9,5,11\\
 631&1,1,1,1,4,7,2,4,10&1259&1,2,4,6,5,11,1,5,17,12,19
\end{array}
\]
The required trial divisors are among
\(2,3,5,7,11,13,17,19,23,29,31\).
For any \(n\) in the small range there are consecutive chain entries
\(p_i\le n<p_{i+1}\). Then \(n<p_{i+1}<2p_i\le2n\).

For real \(x\ge2\), apply the integer statement to
\(n=\lfloor x\rfloor\). An integer \(p>n\) has \(p\ge n+1>x\), and
\(p<2n\le2x\). For \(1<x<2\), take \(p=2\).
\end{proof}

See \cite[Section 8.1, Theorem 8.7]{nzm} for the statement in this
form. The classical central-binomial-coefficient argument uses the
same valuation interface as Theorem~\ref{nt:C08}.
\chapter{Descent and Sums of Squares}
\label{ch:nt-diophantine}

\section{Primitive right triangles and descent}

\begin{theorem}[Primitive Pythagorean triples]\label{nt:D01}
Let \(a,b,c\) be positive integers with \(a^2+b^2=c^2\) and
\(\gcd(a,b)=1\). After interchanging the legs if necessary so that \(b\)
is even, there are unique integers \(r>s>0\) such that
\[
 a=r^2-s^2,\qquad b=2rs,\qquad c=r^2+s^2,\qquad
 \gcd(r,s)=1,\qquad r\not\equiv s\pmod 2.
\]
Conversely these conditions on \(r,s\) give a primitive triple.
\end{theorem}
\begin{proof}
Both legs cannot be even, and two odd squares sum to \(2\bmod4\),
which is not a square. Thus exactly one leg is even and \(c\) is odd.
Any prime dividing \(a,c\) divides \(b^2\) and hence \(b\), so
\(\gcd(a,c)=1\). The positive integers \((c+a)/2\) and \((c-a)/2\)
are coprime because every common divisor divides their sum \(c\)
and their difference \(a\). Their product is \((b/2)^2\).
Unique prime factorization shows that each is a square, say \(r^2\)
and \(s^2\). This also proves uniqueness. They are coprime and their
sum \(c\) is odd, so \(r,s\) have opposite parity.
Conversely the displayed formulas give the identity by expansion.
A common odd prime divisor of \(r^2-s^2\) and \(2rs\) would divide
one of \(r,s\) and then both. The first leg is odd, so the two legs
are coprime.
\end{proof}

\begin{theorem}[Fermat's descent at exponent four]\label{nt:D02}
If \(x,y\in\mathbb Z\setminus\{0\}\), no integer \(z\) satisfies
\(x^4+y^4=z^2\). In particular \(x^4+y^4=z^4\) has no solution with
\(xy\ne0\).
\end{theorem}
\begin{proof}
If a solution exists, take positive \(x,y,z\) by taking absolute
values, and choose one with \(z\) least. If \(d=\gcd(x,y)>1\), then
\(d^4\mid z^2\), whence \(d^2\mid z\) by prime valuations. Dividing
\(x,y\) by \(d\) and \(z\) by \(d^2\) gives a smaller solution.
Thus \((x^2,y^2,z)\) is primitive, and we may take \(x\) odd and
\(y\) even. Theorem~\ref{nt:D01} gives
\[
 x^2=r^2-s^2,\qquad y^2=2rs,\qquad z=r^2+s^2
\]
with \(r,s\) coprime and of opposite parity. If \(r\) were even,
\(r^2-s^2\equiv3\bmod4\), so \(r\) is odd and \(s\) is even.
The identity \((y/2)^2=r(s/2)\) and coprimality give
\(r=u^2\) and \(s=2v^2\), for positive \(u,v\).

The triple \((x,s,r)\) is also primitive: a common prime of \(x,s\)
would divide \(r\) from \(x^2+s^2=r^2\), contrary to \(\gcd(r,s)=1\).
Applying Theorem~\ref{nt:D01} again gives
\[
 x=k^2-h^2,\qquad s=2kh,\qquad r=k^2+h^2,\qquad \gcd(k,h)=1.
\]
Now \(kh=v^2\), so \(k=A^2\) and \(h=B^2\) for positive \(A,B\).
This gives \(A^4+B^4=u^2\) and \(u\le r<z\), contradicting the choice
of \(z\). In the fourth-power equation one can take \(z^2\) as the
square root of its right-hand side and apply the stronger result.
\end{proof}
For the classical parametrization and descent see \cite[Sections
5.3--5.4, Theorems 5.5 and 5.10]{nzm}.

\section{Gaussian integers and two squares}

\begin{theorem}[Arithmetic of Gaussian integers]\label{nt:D03}
The ring \(\mathbb Z[i]\) is Euclidean for the norm
\(N(a+bi)=a^2+b^2\). Its units are \(1,-1,i,-i\), and it has unique
factorization up to units and order. The rational prime \(2\) is a
unit times \((1+i)^2\). An odd rational prime \(p\equiv3\bmod4\)
is prime in \(\mathbb Z[i]\), and if \(p\equiv1\bmod4\) it is the
product of two nonassociate conjugate primes of norm \(p\).
\end{theorem}
\begin{proof}
Complex conjugation gives \(N(zw)=N(z)N(w)\). If \(w\ne0\), round the
two coordinates of \(z/w\) to nearest integers to obtain \(q\in\mathbb
Z[i]\). Then \(z=qw+r\), where
\[
 N(r)=N(w)\,|z/w-q|^2\le\tfrac12N(w)<N(w).
\]
An element is a unit exactly when its norm is \(1\), which lists the
four units. Euclidean division shows that a nonzero ideal is generated
by an element in it of least positive norm. Thus a gcd generates the
ideal of linear combinations and satisfies a B\'ezout identity.
If an irreducible \(\varpi\) does not divide \(a\), their gcd is a unit.
Multiplying a B\'ezout identity by \(b\) then gives
\(\varpi\mid b\) whenever \(\varpi\mid ab\). So irreducibles are
prime. Factoring nonunits by induction on their norms gives existence
and prime cancellation gives uniqueness.

The equation \(2=-i(1+i)^2\) gives the factorization at \(2\), and
\(1+i\) is irreducible because its norm is \(2\). If an odd prime \(p\)
factors into two nonunits, their positive integer norms are both \(p\)
since \(N(p)=p^2\). When \(p\equiv3\bmod4\) this is impossible because
squares modulo \(4\) are \(0,1\). When \(p\equiv1\bmod4\), the
supplement for \(-1\) in Theorem~\ref{nt:B07} gives an integer \(t\)
with \(t^2\equiv-1\bmod p\). The element \(p\) divides
\((t+i)(t-i)\) in \(\mathbb Z[i]\) but divides neither factor, as is
seen from the imaginary coordinates. Therefore \(p\) is reducible.
A factor \(\varpi\) of norm \(p\) satisfies \(p=\varpi\overline\varpi\).
Both factors are irreducible, since a factorization of either into
two nonunits would express its prime norm as a product of integers
greater than \(1\).
If \(\varpi=a+bi\) is associate to its conjugate, a check of the four
units gives \(a=0\), \(b=0\), or \(a=\pm b\). In each case \(p\)
would be a square or twice a square. Thus the factors are nonassociate.
Finally, any Gaussian prime \(\varpi\) divides the positive rational
integer \(N(\varpi)\), so primality and ordinary factorization imply
that it divides a rational prime. The cases above therefore list all
Gaussian prime types.
\end{proof}

\begin{theorem}[Two squares, with representation count]\label{nt:D04}
For \(n>0\), the number \(r_2(n)\) of ordered signed pairs
\((a,b)\in\mathbb Z^2\) with \(a^2+b^2=n\) is
\[
 r_2(n)=4\sum_{d\mid n}\chi_4(d),\qquad
 \chi_4(d)=
 \begin{cases}
 0&d\text{ even},\\
 1&d\equiv1\pmod4,\\
 -1&d\equiv3\pmod4.
 \end{cases}
\]
In particular \(n\) is a sum of two squares exactly when every prime
\(p\equiv3\bmod4\) has even exponent in \(n\). An odd prime is a sum
of two squares exactly when it is \(1\bmod4\). There is exactly one
pair at \(n=0\).
\end{theorem}
\begin{proof}
A pair amounts to a Gaussian integer \(z=a+bi\) of norm \(n\).
In the factorization of \(z\), the exponent of \(1+i\) is fixed
by the exponent of \(2\) in \(n\). An inert rational prime
\(p\equiv3\bmod4\) has norm \(p^2\), so its exponent in \(z\) is
half its exponent in \(n\) and requires that exponent to be even.
A split prime of exponent \(e\) in \(n\) allows \(e+1\) choices
for the exponents of the two norm-\(p\) primes in \(z\). Theorem
\ref{nt:D03} gives all choices independently, and each product gives
four different elements by multiplication by the units, since it is
nonzero. On the other hand \(\chi_4\) is multiplicative: an even
factor gives zero, and the two odd classes multiply as \(1,-1\).
The factor of the divisor sum at \(2\) is \(1\), at a split prime is
\(e+1\), and at an inert prime is \(1-1+\cdots+(-1)^e\).
The factors match the count. The claims about existence, primes and
zero follow.
\end{proof}
See \cite[Chapters XII and XX]{hardy-wright} and
\cite[Section 3.6, Theorems 3.20--3.22 and Corollary 3.23]{nzm}.
The count treats the coordinates as ordered and does not identify
sign changes.

\section{Four squares by a minimal-multiple argument}

For \(A=(a_0,\mathbf a)\) and \(B=(b_0,\mathbf b)\) with
\(\mathbf a,\mathbf b\in\mathbb Z^3\), set
\[
 C(A,B)=
 \bigl(a_0b_0+\mathbf a\mathbin{\cdot}\mathbf b,\ 
       b_0\mathbf a-a_0\mathbf b-\mathbf a\times\mathbf b\bigr).
\]
Write \(\|A\|^2=a_0^2+\|\mathbf a\|^2\).
The cross product is perpendicular to each vector and satisfies
\(\|\mathbf a\times\mathbf b\|^2=
\|\mathbf a\|^2\|\mathbf b\|^2-(\mathbf a\cdot\mathbf b)^2\)
by the coordinate formula. On expanding \(\|C(A,B)\|^2\) the two
terms involving \(2a_0b_0(\mathbf a\cdot\mathbf b)\) cancel,
leaving
\begin{equation}\label{eq:nt-four-square-product}
 \|C(A,B)\|^2=\|A\|^2\|B\|^2.
\end{equation}
In particular all four coordinates of \(C(A,B)\) are integers.
When \(A\equiv B\pmod m\) coordinatewise, its last three coordinates
are \(0\bmod m\), and its first is \(\|B\|^2\bmod m\).

\begin{theorem}[Lagrange's four-square theorem]\label{nt:D05}
Every nonnegative integer is a sum of four squares of nonnegative
integers.
\end{theorem}
\begin{proof}
Zero, one, and the prime \(2\) have explicit representations.
By \eqref{eq:nt-four-square-product} and ordinary factorization
it suffices to handle an odd prime \(p\). The two subsets
\(\{a^2:a\in\mathbb F_p\}\) and \(\{-1-b^2:b\in\mathbb F_p\}\)
of \(\mathbb F_p\) each have \((p+1)/2\) elements, since a nonzero
square has exactly the two roots \(a,-a\). The subsets intersect.
Choose integer representatives \(a,b\) with
\(|a|,|b|\le(p-1)/2\) for an intersection. Then
\[
 a^2+b^2+1=mp,\qquad 0<mp\le 2((p-1)/2)^2+1<p^2.
\]
Thus at least one \(m\) with \(0<m<p\) works. Choose the least
positive \(m\) for which \(mp=\|A\|^2\) with integer coordinates.

If \(m\) is even, the number of odd coordinates of \(A\) is even,
so after reordering they form two pairs of matching parity. A pair
\((a,b)\) is replaced by \(((a+b)/2,(a-b)/2)\), which halves its
sum of squares. This gives a representation of \((m/2)p\),
contrary to minimality. If \(m>1\) is odd, choose \(B\equiv A\bmod m\)
with all coordinates bounded by \((m-1)/2\) in absolute value.
Then \(\|B\|^2=mm'\) with \(0\le m'<m\). If \(m'=0\) then \(m\)
divides every coordinate of \(A\), giving \(m^2\mid mp\) and hence
\(m\mid p\), impossible for \(1<m<p\). Consequently \(0<m'<m\).
All four coordinates of \(C(A,B)\) are divisible by \(m\), by the
congruences after \eqref{eq:nt-four-square-product}. Dividing them
by \(m\) represents \(pm'\), another contradiction. Thus \(m=1\).
Taking absolute values of the final coordinates makes the witnesses
nonnegative.
\end{proof}
For the descent see \cite[Chapter XX]{hardy-wright}.

\section{Primes in reduced residue classes}

We prove the prime-existence result used to choose the auxiliary
prime in Theorem~\ref{nt:D06}. The complex-analytic tools are Taylor
expansion on a disk and holomorphy under locally uniform limits, proved
in the MathGenome complex analysis notes, D12.05.01 and D12.11.02.
The continuation, nonvanishing, and prime-class extraction needed here
are proved below.

\begin{lemma}[Characters of a finite abelian group]
\label{lem:nt-dirichlet-characters}
Let \(G\) be a finite abelian group of order \(h\), and let
\(\widehat G\) be the group of homomorphisms from \(G\) to the
complex numbers of absolute value \(1\). Then \(|\widehat G|=h\),
and
\[
 \frac1h\sum_{\chi\in\widehat G}\chi(g)
  =\begin{cases}1&g=1,\\0&g\ne1.\end{cases}
\]
A nontrivial character satisfies \(\sum_{g\in G}\chi(g)=0\).
If \(g\) has order \(f\), the values \(\chi(g)\) run through all
\(f\)-th roots of unity, each with multiplicity \(h/f\).
\end{lemma}
\begin{proof}
Let \(H\) be a subgroup and \(\psi\) a character of \(H\).
Adjoin \(x\in G\) and let \(d\) be the order of \(xH\) in
\(\langle H,x\rangle/H\). For each of the \(d\) roots \(z\) of
\(z^d=\psi(x^d)\), the formula
\(\widetilde\psi(yx^j)=\psi(y)z^j\) for \(y\in H\) defines a
character on \(\langle H,x\rangle\).
Indeed, \(yx^j=y'x^{j'}\) implies \(j-j'=dk\) for some \(k\)
and \(y'=yx^{dk}\), so both representations give the same value.
Multiplicativity follows because \(G\) is abelian. Every extension
has this form since its value at \(x\) determines it. Iterating
the construction along a chain of subgroups obtained by adjoining
elements gives exactly \([G:H]\) extensions of each character of
\(H\). In particular there are \(h\) characters starting from the
trivial subgroup. Starting from \(\langle g\rangle\) proves the
claim about the distribution of values at \(g\).

For \(g\ne1\), some character \(\chi_*\) has \(\chi_*(g)\ne1\)
by this distribution. Multiplication by \(\chi_*\) permutes the
characters, so the sum of their values at \(g\) equals itself
times \(\chi_*(g)\) and is zero. For \(g=1\) every summand equals
\(1\). Finally, for a nontrivial \(\chi\), choose \(x\) with
\(\chi(x)\ne1\) and reindex the group sum by \(g\mapsto xg\).
The same argument gives \(\sum_{g\in G}\chi(g)=0\).
\end{proof}

\begin{lemma}[Continuation from nonnegative coefficients]
\label{lem:nt-dirichlet-positive-series}
Suppose \(b_n\ge0\) and the series
\(F(s)=\sum_{n\ge1}b_n n^{-s}\) converges absolutely for
\(\operatorname{Re}s>1\), where \(n^{-s}=\exp(-s\log n)\).
If \(F\) has a holomorphic continuation throughout
\(\operatorname{Re}s>0\), the series converges at every positive
real \(\sigma\).
\end{lemma}
\begin{proof}
At \(s=2\) the derivatives are given term by term. For each
derivative order \(j\), the terms on \(\operatorname{Re}s\ge3/2\)
are bounded in absolute value by \(C_j b_n n^{-5/4}\), since
\((\log n)^j n^{-1/4}\) is bounded for \(n\ge1\). This summable
bound justifies the derivative formula on a neighborhood of \(2\).
The disk \(|s-2|<2\) lies in the half-plane of holomorphy,
so its Taylor series represents \(F\) there. For \(0<\sigma\le1\)
put \(t=2-\sigma\), which satisfies \(0<t<2\). The Taylor formula
then gives
\[
 F(\sigma)
 =\sum_{j\ge0}\frac{t^j}{j!}\sum_{n\ge1}b_n(\log n)^j n^{-2}
 =\sum_{n\ge1}b_n n^{-2}\exp(t\log n)
 =\sum_{n\ge1}b_n n^{-\sigma}.
\]
In the first sum the two signs, one from differentiation and one
from \(s-2=-t\), cancel. All terms in the double sum are nonnegative,
so both orders give the same supremum of finite subsums. The
Taylor series converges to a finite value, proving the claim.
For \(\sigma>1\) convergence is an assumption.
\end{proof}

\begin{theorem}[Dirichlet's theorem on primes in arithmetic progressions]
\label{nt:dirichlet-primes}
If \(M\ge1\), \(a\in\mathbb Z\), and \(\gcd(a,M)=1\), then for every
real bound \(T\) there is a prime \(q>T\) with \(q\equiv a\pmod M\).
\end{theorem}
\begin{proof}
Take \(G=(\mathbb Z/M\mathbb Z)^\times\), with the trivial group
for \(M=1\), and let \(h=|G|\). Extend each \(\chi\in\widehat G\)
periodically to integers by using its value on the residue class for
integers coprime to \(M\) and \(0\) otherwise. This extension is
multiplicative since a product is coprime to \(M\) exactly when
both factors are. Write \(\chi_0\) for the principal character,
which is \(1\) on units. For \(\operatorname{Re}s>1\) put
\[
 L(s,\chi)=\sum_{n\ge1}\chi(n)n^{-s}
          =\prod_p(1-\chi(p)p^{-s})^{-1}.
\]
To see the equality, expand the geometric series in the factors
for primes \(p\le Y\). Unique factorization gives the sum over
integers whose prime divisors are all at most \(Y\). The remaining
terms have total absolute value at most
\(\sum_{n>Y}n^{-\operatorname{Re}s}\), which tends to zero.

For a nonprincipal character, the sum over each period is zero
by Lemma~\ref{lem:nt-dirichlet-characters}. Consequently
\(A_\chi(x)=\sum_{n\le x}\chi(n)\) is bounded in absolute value
by \(M\). Summing the integrals of \(s x^{-s-1}\) on \([n,N]\)
over \(1\le n\le N\) gives
\[
 \sum_{n\le N}\chi(n)n^{-s}
   =A_\chi(N)N^{-s}+s\int_1^N A_\chi(x)x^{-s-1}\,dx .
\]
The first term tends to zero on \(\operatorname{Re}s>0\).
On any compact subset of that half-plane, take
\(\operatorname{Re}s\ge\delta>0\) and \(|s|\le B\).
The integral tail from \(N\) is bounded by \(BMN^{-\delta}/\delta\).
The integrals up to finite \(N\) are holomorphic in \(s\), by
expanding the exponential locally uniformly on the bounded
interval of integration. Their locally uniform limit gives the
holomorphic continuation of \(L(s,\chi)\).

For the principal character, its Euler product gives
\(L(s,\chi_0)=\zeta(s)\prod_{p\mid M}(1-p^{-s})\) on
\(\operatorname{Re}s>1\). Applying the same integral identity to
the sequence \(1,1,\ldots\), whose partial-sum function is
\(\lfloor x\rfloor\), gives on that half-plane
\[
 \zeta(s)=s\int_1^\infty\lfloor x\rfloor x^{-s-1}\,dx
         =\frac{s}{s-1}
            -s\int_1^\infty (x-\lfloor x\rfloor)x^{-s-1}\,dx.
\]
The last integral has the locally uniform tail estimate just
proved, with \(1\) in place of \(M\). Thus \(L(s,\chi_0)\) is
meromorphic on \(\operatorname{Re}s>0\) with its only pole there
a simple pole at \(1\). Its residue is
\(c_M=\prod_{p\mid M}(1-p^{-1})>0\), with the empty product
equal to \(1\).

We next prove \(L(1,\chi)\ne0\) for every \(\chi\ne\chi_0\).
Put \(F(s)=\prod_{\chi\in\widehat G}L(s,\chi)\).
At \(p\mid M\) its local factor is \(1\). For \(p\nmid M\) let
\(f_p\) be the order of its class in \(G\). By
Lemma~\ref{lem:nt-dirichlet-characters}, its local factor is
\[
 \prod_{\chi\in\widehat G}(1-\chi(p)p^{-s})^{-1}
     =(1-p^{-f_ps})^{-h/f_p}.
\]
Here the identity follows by multiplying \(1-\omega z\) over
the \(f_p\)-th roots of unity to get \(1-z^{f_p}\).
The resulting power series in \(p^{-s}\) has nonnegative
coefficients, since it is a product of \(h/f_p\) geometric
series in \(p^{-f_ps}\). Its coefficient at every exponent
divisible by \(f_p\) is at least \(1\). Expanding over finite
sets of primes therefore gives multiplicative coefficients
\(b_n\ge0\), with \(b_{m^h}\ge1\) for every \(m\) coprime to \(M\),
since \(f_p\mid h\).

For real \(\sigma>1\), the finite product of local factors
over \(p\le Y\) is bounded by
\(\prod_{p\le Y}(1-p^{-\sigma})^{-h}\le\zeta(\sigma)^h\).
Its expansion is the nonnegative sum over \(n\) with all
prime factors at most \(Y\). Letting \(Y\to\infty\) shows
\(\sum_n b_n n^{-\sigma}<\infty\). This also proves absolute
convergence for complex \(s\) with \(\operatorname{Re}s>1\).
By the Euler products for \(L\) and absolute convergence of the
expanded sum we have \(F(s)=\sum_n b_n n^{-s}\) there.
If a nonprincipal \(L(1,\chi)\) vanished, its local Taylor
expansion would cancel the simple principal pole and make
\(F\) holomorphic throughout \(\operatorname{Re}s>0\).
Lemma~\ref{lem:nt-dirichlet-positive-series} would then give
convergence at \(s=1/h\). But
\[
 \sum_{n\ge1}b_n n^{-1/h}
 \ \ge\ \sum_{\gcd(m,M)=1}\frac{1}{m}
 \ \ge\ \sum_{j\ge0}\frac{1}{1+jM}
 \ =\ \infty.
\]
The last series diverges by comparison with
\(M^{-1}\sum_{j\ge0}(j+1)^{-1}\). The product argument applies
to every nonprincipal character.

For real \(\sigma>1\), define the continuous function
\[
 E_\chi(\sigma)
  =\sum_p\sum_{j\ge1}\frac{\chi(p)^j}{j p^{j\sigma}}.
\]
The double series converges absolutely and locally uniformly
on \(\sigma>1\), by comparison with
\(\sum_{n\ge2}n^{-\sigma}/(1-n^{-\sigma})\).
For \(|z|<1\), the series \(H(z)=\sum_{j\ge1}z^j/j\) satisfies
\(H'(z)=1/(1-z)\). Thus \((1-z)e^{H(z)}\) has derivative zero
and value \(1\) at \(0\). Apply this identity at each prime
with \(z=\chi(p)p^{-\sigma}\) and pass from finite products
to the limit to get
\(\exp E_\chi(\sigma)=L(\sigma,\chi)\).
The terms with \(j\ge2\) are uniformly bounded in sum of absolute
values for \(\sigma>1\) by
\(\sum_{n\ge2}1/(n(n-1))\).

For \(\chi_0\), \(E_{\chi_0}(\sigma)\) is real and hence equals
the real logarithm of \(L(\sigma,\chi_0)\).
The residue computation gives
\[
 E_{\chi_0}(\sigma)=\log\frac1{\sigma-1}+O(1)
 \qquad (\sigma\downarrow1).
\]
For \(\chi\ne\chi_0\), let \(w=L(1,\chi)\ne0\) and choose
\(\lambda\in\mathbb C\) with \(e^\lambda=w\).
On a short interval above \(1\) we have
\(|z(\sigma)|<1/2\) for \(z(\sigma)=L(\sigma,\chi)/w-1\).
The function \(\lambda-H(-z(\sigma))\) is bounded and continuous
there, and its exponential is \(L(\sigma,\chi)\) by the identity
for \(H\). The difference between \(E_\chi\) and
this expression is a continuous function with values in
\(2\pi i\mathbb Z\), so it is constant on that interval.
Thus \(E_\chi(\sigma)=O(1)\) as \(\sigma\downarrow1\).
Subtracting the uniformly bounded tails with \(j\ge2\)
gives the same estimates for \(\sum_p\chi(p)p^{-\sigma}\).

For \(p\nmid M\), the orthogonality identity in
Lemma~\ref{lem:nt-dirichlet-characters} applied to the unit
\([p][a]^{-1}\) selects \(p\equiv a\pmod M\).
For \(p\mid M\) every character term vanishes and such a
prime cannot lie in the reduced class. Thus
\[
 h\sum_{p\equiv a\ (M)}p^{-\sigma}
   =\sum_{\chi\in\widehat G}\overline{\chi(a)}
              \sum_p\chi(p)p^{-\sigma}
   =\log\frac1{\sigma-1}+O(1).
\]
The left side is unbounded as \(\sigma\downarrow1\).
A finite set of primes would give a bounded sum, proving the
asserted existence above any bound.
\end{proof}

For the theorem and a different proof, see \cite[Chapter 7]{apostol}.
The product over all characters is also used in
M.~Stoll's \href{https://mathe2.uni-bayreuth.de/stoll/lecture-notes/IntroductoryNumberTheory.pdf}
{\emph{Introductory Number Theory}}, \S\S23.17--23.19.

\section{Three squares and a three-dimensional lattice}

\begin{lemma}[A three-dimensional lattice bound]\label{lem:nt-minkowski}
Let \(L=B\mathbb Z^3\) for an invertible real \(3\times3\) matrix \(B\).
If an open ball about \(0\) has volume strictly larger than
\(8|\det B|\), then it contains a nonzero point of \(L\).
\end{lemma}
\begin{proof}
The translates by \(L\) of \(F=B[0,1)^3\) partition \(\mathbb R^3\).
This follows by applying \(B^{-1}\) and taking the coordinatewise
integer and fractional parts. The volume of \(F\) is \(|\det B|\).
Let \(U\) be the concentric ball with half the radius of the given
ball, so \(\operatorname{vol}(U)>|\det B|\). The sets
\((U-\lambda)\cap F\), \(\lambda\in L\), have total volume
\(\operatorname{vol}(U)\), by the partition and translation invariance.
Only finitely many of these sets are nonempty, since \(U,F\) are
bounded and \(B^{-1}\) sends bounded sets to bounded sets. They cannot
be pairwise disjoint in \(F\). Thus \(x+\lambda\) and \(x+\mu\) lie
in \(U\) for some \(x\in F\) and \(\lambda\ne\mu\). Their difference
is a nonzero lattice point and its norm is less than twice the radius
of \(U\).
\end{proof}

\begin{theorem}[Legendre's three-square theorem]\label{nt:D06}
A positive integer \(n\) is the sum of three integer squares if and
only if it is not of the form \(4^a(8b+7)\) for nonnegative integers
\(a,b\). Zero is the sum of three squares.
\end{theorem}
\begin{proof}
The triple \((0,0,0)\) represents zero. Squares are \(0,1\bmod4\),
so if three squares sum to a multiple
of \(4\), all three coordinates are even. Repeated division by \(4\)
reduces the claimed obstruction to \(7\bmod8\). The residues of
squares modulo \(8\) are \(0,1,4\). Three such residues cannot sum to
\(7\bmod8\), proving necessity.

For sufficiency, write \(n=d^2 n_0\) with \(n_0\) squarefree.
After removing the powers of \(4\) from \(n\), the remaining
odd square factor is \(1\bmod8\). Thus the obstruction is unchanged
on passing to \(n_0\), and a representation of \(n_0\) gives one of
\(n\) by multiplication by \(d\). We now write \(n\) for the
squarefree integer, with \(n>1\) since \(1\) has an immediate
representation. Write \(n=2^k m\) with \(k\in\{0,1\}\) and \(m\) odd.
Let \(\ell=1\) for \(n\equiv3\bmod8\), and \(\ell=0\) otherwise.

For each odd prime \(p\mid m\) choose a nonzero class of \(q\bmod p\)
such that \(\bigl(\frac{-2^\ell q}{p}\bigr)=1\). For instance one
may impose \(-2^\ell q\equiv1\bmod p\). Also impose \(q\equiv1\bmod8\)
when \(n\equiv1,2,3,5\bmod8\) and \(q\equiv5\bmod8\) when
\(n\equiv6\bmod8\). These are all squarefree possibilities outside
the obstruction. The Chinese remainder theorem gives one class
modulo \(8m\), coprime to \(8m\). Theorem~\ref{nt:dirichlet-primes}
gives a prime \(q>n\) in this class. In particular \(q\equiv1\bmod4\).
For clarity, the relevant signs are recorded together:
\[
\begin{array}{c|ccccc}
 n\bmod 8 &1&2&3&5&6\\ \hline
 k &0&1&0&0&1\\
 (m-1)/2\bmod2 &0&0&1&0&1\\
 \ell &0&0&1&0&0\\
 q\bmod8 &1&1&1&1&5
\end{array}
\]
Hence \(\bigl(\frac{-2^k}{q}\bigr)=(-1)^{k(m-1)/2}\) by the
supplements in Theorem~\ref{nt:B07}. The last sign equals
\(\bigl(\frac{-2^\ell}{m}\bigr)_{\!J}\), with the Jacobi symbol at
\(m=1\) read as \(1\). Indeed, for \(\ell=0\) this is the supplement
for \(-1\); for \(\ell=1\) one has \(m=n\equiv3\bmod8\), and the
two supplements give \(\bigl(\frac{-2}{m}\bigr)_{\!J}=1\).
Reciprocity with \(q\equiv1\bmod4\) and
\(\bigl(\frac{-2^\ell q}{p}\bigr)=1\) give
\[
 \left(\frac{m}{q}\right)
  =\prod_{p\mid m}\left(\frac{p}{q}\right)
  =\prod_{p\mid m}\left(\frac{q}{p}\right)
  =\left(\frac{-2^\ell}{m}\right)_{\!J}.
\]
Thus \(\bigl(\frac{-n}{q}\bigr)=1\).

Set \(b_0=2^\ell q\), which is coprime to \(n\). There are integers
\(u,v\) such that
\begin{equation}\label{eq:nt-three-square-congruences}
 u^2\equiv-n\pmod{4^\ell q},\qquad
 v^2\equiv-b_0^{-1}\pmod n .
\end{equation}
For the first congruence, a root exists modulo \(q\) by the last
symbol. When \(\ell=1\), \(-n\equiv1\bmod4\) and we combine a root
modulo \(4\) with the root modulo \(q\). For the second congruence,
at each odd \(p\mid n\) the element \(-b_0^{-1}\) is a square
because \(-b_0\) is a nonzero square there. If \(2\mid n\), then
\(b_0\) is odd and the congruence is \(v^2\equiv1\bmod2\).
Squarefreeness and CRT combine the roots modulo all of \(n\).

For integers \(r,s,t\) put \(W=b_0r+us\) and consider the lattice
of triples
\[
 X=vW+nt,\qquad Y=\frac{W}{\sqrt{b_0}},\qquad
 Z=\sqrt{\frac n{b_0}}\,s .
\]
The linear transformation has determinant of absolute value
\(n\sqrt{b_0}\sqrt{n/b_0}=n^{3/2}\), by expansion along its last
column. An open ball of radius \(\sqrt{2n}\) has volume
\(\frac{8\sqrt2\pi}{3}n^{3/2}>8n^{3/2}\). The formula follows by
integrating the area \(\pi(2n-z^2)\) of its horizontal disk
sections, and the strict inequality follows from \(\pi>3\).
Lemma~\ref{lem:nt-minkowski} gives a nonzero lattice point in this
ball. At this point \(X\) is an integer and
\[
 S=Y^2+Z^2
  =\frac{W^2+ns^2}{b_0}
  =b_0r^2+2urs+\frac{u^2+n}{b_0}s^2
\]
is an integer by \eqref{eq:nt-three-square-congruences}.
Modulo \(n\), one has \(S\equiv b_0^{-1}W^2\) and
\(X\equiv vW\), so \(X^2+S\equiv0\bmod n\).
The squared norm is strictly between \(0\) and \(2n\);
consequently \(X^2+S=n\).

It remains to prove \(S\) is a sum of two integer squares. If
\(S=0\) this holds. Otherwise let \(p\) be an odd prime of odd
exponent in \(S\). Since \(q>n\ge S\), we have \(p\nmid b_0\).
First suppose \(p\nmid n\). From \(n=X^2+S\) it follows that
\(\bigl(\frac n p\bigr)=1\). Remove the largest common power \(p^j\)
from \(W,s\), where \(j\) exists since \(S>0\). Write
\(W=p^j W_1,\ s=p^j s_1\). The exponent of \(p\) in
\(W_1^2+ns_1^2\) is still odd, since the original exponent in
\(b_0S\) was odd and we removed \(p^{2j}\). In particular
\(W_1^2\equiv-ns_1^2\bmod p\). Neither coordinate is zero modulo
\(p\), since otherwise both would be, contrary to the choice of \(j\).
Thus \(\bigl(\frac{-n}{p}\bigr)=1\), and division by the first
symbol yields \(\bigl(\frac{-1}{p}\bigr)=1\).

Now suppose \(p\mid n\). Since \(p\mid S\), reduction modulo \(p\) gives \(p\mid X\)
and \(p\mid W\) from \(X^2+S=n\) and \(b_0 S=W^2+ns^2\).
Substituting the latter into \(b_0n=b_0X^2+b_0S\) gives
\(b_0n\equiv ns^2\bmod p^2\). Because \(p\) divides the squarefree
\(n\) to exponent one, division yields \(s^2\equiv b_0\bmod p\).
So \(\bigl(\frac{b_0}{p}\bigr)=1\), and the imposed condition
\(\bigl(\frac{-b_0}{p}\bigr)=1\) again gives
\(\bigl(\frac{-1}{p}\bigr)=1\).
In both cases the supplement gives \(p\equiv1\bmod4\). Thus the
criterion in Theorem~\ref{nt:D04} applies to \(S\). Adding \(X^2\)
gives the desired three squares, and the square factor removed at
the start can be restored.
\end{proof}
The sufficiency argument is the lattice proof of Ankeny
\cite{ankeny}, in the uniform modular setup of \cite[Section 5]{fan}.
The inequalities, congruences, and final use of the two-square
criterion above also cover the even squarefree cases.
\chapter[Continued fractions and approximation]{Continued fractions, approximation, and irrationality}
\label{ch:nt-cf}

We write a rational number as \(p/q\) with \(p,q\in\mathbb Z\) and
\(q>0\). Its reduced form satisfies \(\gcd(p,q)=1\). The absolute errors
\(\lvert x-p/q\rvert\) and \(\lvert qx-p\rvert\) lead to different
approximation statements. We keep the two errors separate throughout.

\section{Rational roots and classical irrational numbers}
\label{nt:E01}

\begin{theorem}[Rational roots]
Let \(P(X)=c_dX^d+\cdots+c_0\in\mathbb Z[X]\), where \(d\ge1\) and
\(c_d\ne0\). If \(p/q\) is a nonzero reduced root with \(q>0\), then
\(p\mid c_0\) and \(q\mid c_d\). If \(0\) is a root, then \(c_0=0\) and
its reduced denominator is \(1\). In particular, every rational root of
a monic integer polynomial is an integer.
\end{theorem}
\begin{proof}
Multiplication of \(P(p/q)=0\) by \(q^d\) gives
\[
 c_dp^d+c_{d-1}p^{d-1}q+\cdots+c_0q^d=0.
\]
Reduction modulo \(\lvert p\rvert\) gives \(p\mid c_0q^d\), and reduction
modulo \(q\) gives \(q\mid c_dp^d\). Coprimality and Bézout's identity
permit cancellation of \(q^d\) in the first divisibility and \(p^d\) in
the second. The zero and monic cases follow from evaluation at zero
and the denominator condition.
\end{proof}

\begin{corollary}
If \(k\ge2\) and \(m\) is a nonnegative integer which is not a \(k\)th
power in \(\mathbb Z_{\ge0}\), then the nonnegative real \(k\)th root of
\(m\) is irrational. This includes \(\sqrt m\) for nonsquare \(m\).
\end{corollary}
\begin{proof}
A rational nonnegative root of \(X^k-m\) would be a nonnegative integer
by the monic case. Its \(k\)th power would equal \(m\).
\end{proof}

\begin{theorem}
The real numbers \(e\) and \(\pi\) are irrational.
\end{theorem}
\begin{proof}
The Taylor series for \(\exp(1)\) gives
\(e=\sum_{j=0}^{\infty}1/j!\). If \(e=a/b\) with \(b>0\), choose an
integer \(N\ge\max\{b,2\}\). Both \(N!e\) and
\(N!\sum_{j=0}^N1/j!\) would be integers, whereas their difference
satisfies
\[
 0<\sum_{k=1}^{\infty}\frac{1}{(N+1)\cdots(N+k)}
 \le \sum_{k=1}^{\infty}(N+1)^{-k}=\frac1N<1.
\]

For \(\pi\), suppose \(\pi=a/b\) with positive integers \(a,b\). For
a positive integer \(N\) set
\[
 f(t)=\frac{t^N(a-bt)^N}{N!},\qquad
 F(t)=\sum_{j=0}^{N}(-1)^j f^{(2j)}(t).
\]
At \(t=0\) every derivative of \(f\) is an integer. In fact the
coefficient of \(t^{N+k}\), for \(0\le k\le N\), is an integer divided
by \(N!\), and \((N+k)!/N!\) is an integer. The remaining derivatives
vanish at zero. Since \(f(\pi-t)=f(t)\), the endpoint derivatives at
\(\pi\) are integers as well. The degree of \(f\) is \(2N\), so
\(F''+F=f\). Consequently
\[
 \int_0^\pi f(t)\sin t\,dt
 =\bigl[F'(t)\sin t-F(t)\cos t\bigr]_0^\pi
 = F(\pi)+F(0)\in\mathbb Z.
\]
On \(0<t<\pi\) the integrand is strictly positive. The bound
\(t(a-bt)=bt(\pi-t)\le b\pi^2/4\) gives
\[
 0<\int_0^\pi f(t)\sin t\,dt
 \le \frac{\pi(b\pi^2/4)^N}{N!}.
\]
The successive upper bounds have ratio \((b\pi^2/4)/(N+1)\), which
tends to zero. For large \(N\) the integral is an integer strictly
between \(0\) and \(1\), a contradiction.
\end{proof}

For the rational-root theorem and the proofs of irrationality, see
\cite[\S6.3, Theorems 6.13 and 6.18 and Exercise 8]{nzm}.

\section{Finite simple continued fractions}
\label{nt:E02}

For a real number \(x\), set \(\alpha_0=x\). At step \(j\), put
\(a_j=\lfloor\alpha_j\rfloor\). If \(\alpha_j=a_j\), stop. Otherwise
set
\[
 \alpha_{j+1}=\frac{1}{\alpha_j-a_j}.
\]
The \(\alpha_j\) are the complete quotients and the \(a_j\) the partial
quotients. Whenever it is defined, \(\alpha_{j+1}>1\), so \(a_j\ge1\)
for \(j\ge1\). The initial integer \(a_0\) can have any sign. For a finite
list we use the notation
\[
 [a_0;a_1,\ldots,a_n]
 =a_0+\cfrac1{a_1+\cfrac1{\ddots+\cfrac1{a_n}}},
 \qquad [a_0]=a_0.
\]

\begin{theorem}[Finite expansion and normalization]
The algorithm terminates exactly for rational \(x\). Every rational
has a unique finite expansion with \(a_0\in\mathbb Z\), \(a_j\ge1\) for
\(j\ge1\), and \(a_n\ge2\) if \(n\ge1\). For \(a_n\ge2\) the
expansion obtained by replacing \(a_n\) by \(a_n-1,1\) has the same
value. For an integer \(a_0\), its other finite expansion is
\([a_0-1;1]\).
\end{theorem}
\begin{proof}
For a rational complete quotient \(u/v\) in lowest terms with \(v>0\),
Euclidean division with remainder gives \(u=a_jv+r\), \(0\le r<v\),
including when \(u<0\). If \(r=0\), the algorithm stops. If \(r>0\),
the next complete quotient is \(v/r\), and \(\gcd(v,r)=1\).
The positive denominators strictly decrease. Conversely, termination
expresses \(x\) as a finite list of additions and nonzero reciprocals of
integers, hence as a rational number. If the algorithm stops at
\(j\ge1\), then \(\alpha_j>1\) is an integer and \(a_j\ge2\).

For a list satisfying the last-entry condition, each tail with index
at least \(1\) is greater than \(1\). Thus each entry is the floor of the
corresponding tail, and reciprocating its fractional part recovers the
next tail. The given list must therefore be exactly the one produced
by the algorithm. The identities for a final \(1\) follow from
\(u+1/1=u+1\). Conversely a noncanonical list ends in \(1\), and merging
its last two entries produces a canonical list, including the case
of a two-entry list that merges to an integer. This accounts for all
finite expansions under the stated digit conventions.
\end{proof}

See \cite[Chapter X]{hardy-wright} and \cite[\S\S7.1--7.3]{nzm}.

\section{Convergents and exact error bounds}
\label{nt:E03}

\begin{theorem}[Convergent recurrences and estimates]
For the simple continued fraction of \(x\), define
\[
 \begin{aligned}
 p_{-2}&=0,&p_{-1}&=1,&p_n&=a_np_{n-1}+p_{n-2},\\
 q_{-2}&=1,&q_{-1}&=0,&q_n&=a_nq_{n-1}+q_{n-2}.
 \end{aligned}
\]
At every defined index \(n\ge0\), the convergent \(p_n/q_n\) is
\([a_0;\ldots,a_n]\) in lowest terms, and
\begin{equation}
 p_nq_{n-1}-p_{n-1}q_n=(-1)^{n+1}.
 \label{eq:nt-cf-det}
\end{equation}
The denominators satisfy \(q_0=1\), \(q_1=a_1\) if defined, and
\(q_{n+1}>q_n\) for \(n\ge1\) when both exist. If \(x\) is irrational,
then for all \(n\ge0\)
\begin{equation}
 x=\frac{p_n\alpha_{n+1}+p_{n-1}}{q_n\alpha_{n+1}+q_{n-1}},\qquad
 q_nx-p_n=\frac{(-1)^n}{q_n\alpha_{n+1}+q_{n-1}},
 \label{eq:nt-cf-tail}
\end{equation}
and
\begin{equation}
 \frac{1}{q_n(q_{n+1}+q_n)}
 < \left|x-\frac{p_n}{q_n}\right|
 < \frac{1}{q_nq_{n+1}}.
 \label{eq:nt-cf-error}
\end{equation}
In particular the convergents tend to \(x\).
\end{theorem}
\begin{proof}
The recurrences give, by induction,
\[
 \prod_{j=0}^{n}
 \begin{pmatrix}a_j&1\\1&0\end{pmatrix}
 =\begin{pmatrix}p_n&p_{n-1}\\q_n&q_{n-1}\end{pmatrix},
\]
where the factors are in increasing order from left to right.
This is the matrix of the composition of the maps \(t\mapsto a_j+1/t\).
Evaluating the finite nested expression gives \(p_n/q_n\).
The determinant of each factor is \(-1\), proving
\eqref{eq:nt-cf-det}. Any common divisor of \(p_n,q_n\) divides this
determinant, so the fraction is reduced. The recurrence gives
\(q_0=1\) and \(q_1=a_1\ge1\). For \(n\ge1\) we have
\(q_{n+1}\ge q_n+q_{n-1}>q_n\).

In the infinite case, evaluating the composition at
\(t=\alpha_{n+1}\) proves the first equality in
\eqref{eq:nt-cf-tail}. Subtracting \(p_n\) after multiplication by
\(q_n\) and using \eqref{eq:nt-cf-det} proves the second.
Now \(a_{n+1}<\alpha_{n+1}<a_{n+1}+1\), whence
\[
 q_{n+1}<q_n\alpha_{n+1}+q_{n-1}<q_{n+1}+q_n.
\]
The reciprocal inequalities yield \eqref{eq:nt-cf-error}. The
denominators grow unboundedly since they are eventually strictly
increasing positive integers.
\end{proof}

In particular, an infinite digit string produced by the algorithm
determines its real value uniquely, by the limit of its convergents.
The recurrences, determinants, and tail formulas can be compared with
\cite[\S\S7.2--7.4, Theorems 7.4, 7.5, and 7.11]{nzm}.

\section{Best approximation and Legendre's criterion}
\label{nt:E04}

\begin{theorem}[Best approximation of the second kind]
Let \(x\) be irrational and \(p_n/q_n\) be its \(n\)th convergent. If
\(p,q\) are integers and \(0<q<q_{n+1}\), then
\[
 \lvert qx-p\rvert\ge \lvert q_nx-p_n\rvert.
\]
Equality is possible only for \((p,q)=(p_n,q_n)\) under the indicated
denominator constraint.
\end{theorem}
\begin{proof}
Consecutive column vectors \((p_n,q_n)\) and \((p_{n+1},q_{n+1})\)
are an integer basis of \(\mathbb Z^2\) because their determinant is
\(1\) or \(-1\). Write
\((p,q)=u(p_n,q_n)+v(p_{n+1},q_{n+1})\) with integers \(u,v\).
Put \(\delta_i=q_ix-p_i\). By \eqref{eq:nt-cf-tail},
\(\delta_n,\delta_{n+1}\) are nonzero with opposite signs.
If \(v=0\), then \(q=uq_n>0\), so \(u\ge1\) and the error is
\(u\lvert\delta_n\rvert\).
If \(v>0\), then \(q<q_{n+1}\) forces \(u<0\). If \(v<0\), then
\(q>0\) forces \(u>0\). In either case \(v\ne0\), the terms
\(u\delta_n\) and \(v\delta_{n+1}\) have the same sign. Their combined
absolute value is strictly greater than
\(\lvert u\rvert\lvert\delta_n\rvert\ge\lvert\delta_n\rvert\).
\end{proof}

\begin{corollary}[Legendre's criterion]
\label{cor:nt-legendre}
If \(x\) is irrational and a reduced \(p/q\) with \(q>0\) satisfies
\[
 \left|x-\frac pq\right|<\frac{1}{2q^2},
\]
then \(p/q\) is a convergent of \(x\).
\end{corollary}
\begin{proof}
Choose the largest \(n\ge0\) with \(q_n\le q\), which exists since
\(q_0=1\) and \(q_n\to\infty\). The choice handles the possible equality
\(q_0=q_1\) and gives \(q<q_{n+1}\). If \(p/q\ne p_n/q_n\), the
nonzero integer \(q_np-qp_n\) satisfies
\[
 \begin{aligned}
 1
 &\le \lvert q_np-qp_n\rvert\\
 &\le q_n\lvert qx-p\rvert+q\lvert q_nx-p_n\rvert\\
 &\le (q_n+q)\lvert qx-p\rvert<1.
 \end{aligned}
\]
The middle estimate uses the second-kind theorem, and the last uses
\(q_n\le q\) and the assumed strict bound. This contradiction gives the
required convergent.
\end{proof}

See \cite[\S7.5, Theorems 7.13--7.14]{nzm}.

\section{Dirichlet's approximation theorem}
\label{nt:E05}

\begin{theorem}[Bounded denominators]
For each real \(x\) and integer \(N\ge1\), there are integers \(p,q\)
with \(1\le q\le N\) and
\[
 \lvert qx-p\rvert\le\frac1{N+1}.
\]
The bound is strict when \(x\) is irrational. Consequently an irrational
\(x\) has infinitely many distinct reduced \(p/q\) with
\(\lvert x-p/q\rvert<1/q^2\).
\end{theorem}
\begin{proof}
Take the \(N+1\) fractional parts \(\{jx\}\) for \(0\le j\le N\),
with labels \(j\). If two coincide, their distinct labels give
\(qx=p\) by subtraction and a sign choice, with \(1\le q\le N\).
Otherwise order the points cyclically on the unit interval with
endpoints identified. The \(N+1\) successive gaps, including the gap
from the last point back to the first, sum to \(1\). A shortest gap
has length at most \(1/(N+1)\). With endpoint labels \(i,j\), that gap
equals \((j-i)x\) minus an integer, with the integer chosen to account
for any wrap. Taking \(q=\lvert j-i\rvert\) and changing the sign of
the integer if needed gives the claimed absolute value. If \(x\) is
irrational, every gap is irrational, since \(j-i\ne0\), and cannot
equal \(1/(N+1)\).

For irrational \(x\), the fractions thus obtained satisfy
\(\lvert x-p/q\rvert\le1/(q(N+1))<1/q^2\). Reducing the fraction
preserves the strict \(1/q^2\) estimate when \(q\) is replaced by its
possibly smaller reduced denominator. As \(N\to\infty\), the errors
tend to zero because they are at most \(1/(N+1)\). A finite collection
of rationals would have a positive minimum distance from \(x\), so
infinitely many different reduced fractions occur.
\end{proof}

Compare \cite[\S6.2, Theorems 6.8--6.9]{nzm}.

\section{Farey fractions}
\label{nt:E06}

Let \(\mathcal F_N\) be the increasing list of reduced fractions in
\([0,1]\) with denominator at most \(N\), where \(N\ge1\).

\begin{theorem}[Neighbors, mediants, and insertion]
If \(a/b<c/d\) are consecutive members of \(\mathcal F_N\), then
\[
 bc-ad=1,\qquad b+d>N.
\]
The unique reduced fraction strictly between them of least denominator
is the mediant \((a+c)/(b+d)\). To obtain \(\mathcal F_{N+1}\), insert
this mediant in precisely the neighboring intervals with \(b+d=N+1\).
\end{theorem}
\begin{proof}
First suppose \(a/b<c/d\) are reduced with \(bc-ad=1\). For any reduced
\(p/q\) strictly between them, the vectors \((a,b),(c,d)\) form an
integer basis, and
\[
 (p,q)=u(a,b)+v(c,d),\qquad
 u=cq-dp>0,\qquad v=bp-aq>0.
\]
The coefficients are positive integers because \(a/b<p/q<c/d\).
It follows that \(q=ub+vd\ge b+d\), with equality exactly for
\(u=v=1\). The mediant is reduced, since a common divisor of
\(a+c,b+d\) divides \((a+c)b-(b+d)a=1\). The determinant of each of
the pairs formed by splitting at the mediant is also \(1\), as direct
substitution shows.

At \(N=1\), the endpoints \(0/1,1/1\) have determinant \(1\) and
denominator sum \(2>1\). Assume the claims about neighbors at \(N\).
Any fraction new at \(N+1\) has denominator \(N+1\) and lies in the
interior of one of these neighboring intervals. The first paragraph
shows that such a fraction exists in this interval exactly when
\(b+d=N+1\) and that it then is the mediant. Pairs untouched by
insertion have sum \(b+d>N+1\). Pairs created by insertion have sums
\(2b+d\) and \(b+2d\), both larger than \(b+d=N+1\), and determinant
\(1\). This proves the result by induction.
\end{proof}

See \cite[Chapter III]{hardy-wright} and
\cite[\S\S6.1--6.2, Theorems 6.1, 6.4, and 6.7]{nzm}.

\section{Hurwitz's bound and its optimality}
\label{nt:E07}

\begin{theorem}[Hurwitz]
For every irrational real \(x\), infinitely many distinct reduced
fractions \(p/q\) satisfy
\[
 \left|x-\frac pq\right|<\frac{1}{\sqrt5\,q^2}.
\]
For any \(C>\sqrt5\), the assertion with \(C\) in place of \(\sqrt5\)
fails for \(x=(1+\sqrt5)/2\).
\end{theorem}
\begin{proof}
First subtract an integer from \(x\) so that \(0<x<1\). This does not
alter denominators, coprimality, or approximation errors after
translating fractions back. For each \(N\), choose the neighbors
\(a/b<c/d\) of \(\mathcal F_N\) bracketing \(x\). Their distance is
\(1/(bd)\). Put \(g=(1+\sqrt5)/2\) and \(s=d/b\). If both neighbors
fail the strict estimate, then
\[
 \frac{1}{bd}
 =\left(x-\frac ab\right)+\left(\frac cd-x\right)
 \ge\frac1{\sqrt5}\left(\frac1{b^2}+\frac1{d^2}\right),
 \qquad
 s+\frac1s\le\sqrt5.
\]
The latter inequality is \(s^2-\sqrt5\,s+1\le0\), whose roots are
\(g^{-1}\) and \(g\). Both endpoints are irrational, while \(s\) is
rational, so \(g^{-1}<s<g\).

If \(x\) is below the mediant and the mediant also fails, the left
neighbor and the mediant bracket \(x\). These have determinant \(1\),
so the same necessary inequality applies to the ratio
\((b+d)/b=1+s>1+g^{-1}=g\), a contradiction. If \(x\) is above the
mediant, use the right neighbor and the mediant with ratio
\((b+d)/d=1+1/s>g\). The gap test is invariant under inverting the
ratio, so it again gives a contradiction. The irrational \(x\) cannot
equal the mediant. Hence one of the two endpoints or their mediant is
a successful fraction.

All three candidates lie in a closed interval of length \(1/(bd)\).
From \(b+d\ge N+1\) and \((b-1)(d-1)\ge0\) we get
\(bd\ge b+d-1\ge N\). Thus the successful fractions have errors at
most \(1/N\). These fractions cannot belong to a fixed finite set,
whose distances from \(x\) would have a positive minimum.

For sharpness write \(g'=(1-\sqrt5)/2\). The polynomial
\(T^2-T-1\) has irrational roots \(g,g'\). For each rational \(p/q\),
\[
 1\le \lvert p^2-pq-q^2\rvert
 =q^2\left|\frac pq-g\right|\left|\frac pq-g'\right|.
\]
If infinitely many distinct reduced fractions satisfied the stronger
bound with \(C>\sqrt5\), one could take a sequence with denominators
tending to infinity. Indeed, for bounded denominators there are only
finitely many fractions in the bounded interval imposed by the bound.
Along this sequence \(p/q\to g\), and the displayed inequality would
force \(\lvert p/q-g'\rvert>C\) at every term. Its limit is
\(\lvert g-g'\rvert=\sqrt5<C\), a contradiction.
\end{proof}

This is the Farey-neighbor route of
\cite[\S6.2, Theorems 6.11--6.12]{nzm}, with the gap comparison
expressed through the two roots of \(s^2-\sqrt5\,s+1\).

\section{Periodicity and reduced complete quotients}

A quadratic irrational is a real number of degree two over
\(\mathbb Q\). Its conjugate \(y'\) is the other root of its irreducible
quadratic. We call \(y\) reduced when \(y>1\) and \(-1<y'<0\).
An infinite continued fraction is eventually periodic if its digit
sequence is periodic after finitely many entries, and purely periodic
if periodic from its first entry.

\begin{lemma}[Reduced quotients and their predecessors]
\label{lem:nt-reduced-step}
The map \(T(y)=1/(y-\lfloor y\rfloor)\) sends reduced quadratic
irrationals to reduced quadratic irrationals. Each reduced \(z\) has
exactly one reduced predecessor under \(T\), namely
\[
 y=\left\lfloor-\frac1{z'}\right\rfloor+\frac1z .
\]
If \(y=(P+\sqrt D)/Q\), with integers \(P,Q,D\), positive nonsquare
\(D\), \(Q\ne0\), and \(Q\mid D-P^2\), then its next quotient has
the same \(D\) and is
\begin{equation}
 T(y)=\frac{P_++\sqrt D}{Q_+},
 \qquad
 P_+=\lfloor y\rfloor Q-P,\qquad Q_+=\frac{D-P_+^2}{Q}.
 \label{eq:nt-cf-surd-step}
\end{equation}
Here \(Q_+\ne0\) is an integer and \(Q_+\mid D-P_+^2\).
\end{lemma}
\begin{proof}
For reduced \(y\), we have \(a=\lfloor y\rfloor\ge1\) and
\(0<y-a<1\), so \(T(y)>1\). Its conjugate is \(1/(y'-a)\in(-1,0)\).
For a reduced \(z\), a predecessor with digit \(a\) must be
\(y=a+1/z\). Its conjugate lies in \((-1,0)\) exactly when
\[
 -1-\frac1{z'}<a<-\frac1{z'} .
\]
Since \(-1/z'>1\) is irrational, the open interval contains exactly
the integer \(a=\lfloor-1/z'\rfloor\ge1\). Then \(y>1\) and its floor
is \(a\) because \(0<1/z<1\). This proves existence and uniqueness.

For the formula, \(y-a=(\sqrt D-P_+)/Q\); rationalizing gives
\eqref{eq:nt-cf-surd-step}. The equality \(P_+\equiv-P\bmod Q\)
ensures that \(Q_+\) is an integer. Nonsquareness of \(D\) makes
\(D-P_+^2\ne0\). The relation \(D-P_+^2=QQ_+\) gives the divisibility
in the next state.
\end{proof}

\begin{theorem}[Lagrange periodicity and the pure-period criterion]
\label{nt:E08}
The simple continued fraction of a real irrational \(x\) is
eventually periodic if and only if \(x\) is a quadratic irrational.
For a real quadratic irrational \(x\), the expansion is purely
periodic if and only if \(x\) is reduced.
\end{theorem}
\begin{proof}
First suppose a complete quotient \(y\) has a purely periodic expansion
with period length \(r\ge1\), and let \(b_0,\ldots,b_{r-1}\) be the
first \(r\) digits. Since the first digit recurs at a positive index,
all these digits are positive. In particular \(y>1\). Its complete
quotient at index \(r\) is again \(y\), because their infinite digit
strings agree and their convergents determine the value. Let
\(P_i^\ast,Q_i^\ast\) be the numerator and denominator recurrences for
this digit list, with the initial conventions of
Section~\ref{nt:E03}, and put \(t=r-1\). The fixed-point equality
from \eqref{eq:nt-cf-tail} is equivalent to
\[
 Q_t^\ast y^2+(Q_{t-1}^\ast-P_t^\ast)y-P_{t-1}^\ast=0.
\]
This is a genuine quadratic since \(Q_t^\ast>0\). As \(y\) is
irrational, this polynomial gives its quadratic conjugate.
The polynomial is negative at \(0\) because \(P_{t-1}^\ast>0\).
At \(-1\) its value is
\((Q_t^\ast-Q_{t-1}^\ast)+(P_t^\ast-P_{t-1}^\ast)>0\).
For \(t=0\) the denominator difference is \(1\) and the numerator
difference is \(b_0-1\ge0\). For \(t=1\) the numerator difference is
positive because \(P_1^\ast=b_1b_0+1\). For \(t\ge2\) both differences
are positive by the recurrences. Thus the root other than \(y>1\)
belongs to \((-1,0)\), proving that \(y\) is reduced.
If \(x\) has an eventually periodic expansion, some complete quotient
has a purely periodic expansion. Formula \eqref{eq:nt-cf-tail}
expresses \(x\) as an integer linear fractional transformation of this
quadratic irrational when the periodic part starts after the first
digit. Hence \(x\) lies in a quadratic extension of \(\mathbb Q\);
since it is irrational, it has degree two. When the periodic part
starts at the first digit, the preceding argument applies to \(x\)
directly.

Now let \(x\) be a quadratic irrational. If \(Ax^2+Bx+C=0\)
with integers \(A\ne0,B,C\), its discriminant
\(D=B^2-4AC\) is a positive nonsquare integer. The quadratic formula
writes \(x=(P_0+\sqrt D)/Q_0\) with \(P_0=-B,Q_0=2A\) for the
plus root, or \(P_0=B,Q_0=-2A\) for the minus root. In both cases
\(Q_0\mid D-P_0^2\). The recurrence
\eqref{eq:nt-cf-surd-step} gives this representation for every complete
quotient.

To see that the quotients eventually become reduced, solve
\eqref{eq:nt-cf-tail} for \(\alpha_{n+1}\) and conjugate, obtaining
\[
 \alpha_{n+1}'
   =\frac{p_{n-1}-q_{n-1}x'}{q_nx'-p_n}.
\]
The denominator is nonzero because \(x'\) is irrational.
For large \(n\), both \(x'-p_{n-1}/q_{n-1}\) and
\(x'-p_n/q_n\) have the same sign as \(x'-x\), by convergence of the
convergents. Since the denominator indices are then positive, the
displayed quotient is negative. A further step has conjugate
\(1/(\alpha_{n+1}'-a_{n+1})\in(-1,0)\), because \(a_{n+1}\ge1\).
The real quotient at this step is greater than \(1\), so the reduced
condition holds and remains true by Lemma~\ref{lem:nt-reduced-step}.

For any reduced state \((P+\sqrt D)/Q\), the difference of its two
conjugates is \(2\sqrt D/Q>0\), so \(Q>0\). The real quotient greater
than \(1\) gives \(Q<P+\sqrt D\). The negative conjugate gives
\(P<\sqrt D\). In particular \(-\sqrt D<P<\sqrt D\) and
\(0<Q<2\sqrt D\). There are only finitely many such integer pairs.
The update \eqref{eq:nt-cf-surd-step} is deterministic, so the sequence
of reduced states repeats and the digits are eventually periodic.

If \(x\) was already reduced, every complete quotient in its orbit is
reduced. Suppose \(T^i(x)=T^j(x)\) with \(0\le i<j\).
Uniqueness of the reduced predecessor gives
\(T^{i-1}(x)=T^{j-1}(x)\) when \(i>0\). Repeating gives
\(x=T^{j-i}(x)\). Thus its digit sequence was periodic from the first
entry. The necessity of the reduced condition for a pure period was
proved in the first paragraph.
\end{proof}

The two periodicity criteria are given in
\cite[\S7.7, Theorems 7.19--7.20]{nzm}. In the case \(x=\sqrt D\),
our notation gives \(P_0=0,Q_0=1\) in
\eqref{eq:nt-cf-surd-step}. The next chapter uses this recurrence to
identify exactly which convergents solve the Pell equations.

\section{Liouville's inequality and an explicit transcendental number}
\label{nt:E09}

\begin{theorem}[Liouville's inequality]
If \(x\) is a real algebraic irrational of degree \(d\ge2\), there
exists \(c_x>0\) such that for every reduced \(p/q\) with \(q>0\),
\[
 \left|x-\frac pq\right|\ge \frac{c_x}{q^d}.
\]
\end{theorem}
\begin{proof}
Let \(P\) be an irreducible integer polynomial of degree \(d\) with
\(P(x)=0\), obtained by clearing denominators in the minimal
polynomial. No rational is a root, since irreducibility and a
rational root would give degree one. Thus
\(q^dP(p/q)\) is a nonzero integer and
\(\lvert P(p/q)\rvert\ge q^{-d}\). Set
\[
 M=\max\left\{1,\sup_{t\in[x-1,x+1]}\lvert P'(t)\rvert\right\}.
\]
If \(p/q\) lies in that interval, the mean value theorem gives
\(\lvert P(p/q)\rvert\le M\lvert x-p/q\rvert\).
Outside the interval, \(\lvert x-p/q\rvert>1\ge1/(Mq^d)\).
Take \(c_x=1/M\).
\end{proof}

\begin{corollary}
The number \(\displaystyle L=\sum_{j=1}^{\infty}10^{-j!}\) is
transcendental.
\end{corollary}
\begin{proof}
Put \(Q_N=10^{N!}\), and let \(P_N/Q_N\) be the truncation after the
\(N\)th term, with \(P_N\in\mathbb Z\). The positive tail satisfies
\[
 0<L-\frac{P_N}{Q_N}
 \le\frac{10}{9}\,10^{-(N+1)!}
 <2Q_N^{-(N+1)}.
\]
The first inequality follows by bounding the tail by the geometric
series of all integer exponents starting at \((N+1)!\).
If \(L=a/b\) were rational with \(b>0\), its positive difference
from \(P_N/Q_N\) would be at least \(1/(bQ_N)\), contradicting the
upper bound for large \(N\). Thus \(L\) is irrational.
If it were algebraic of degree \(d\), Liouville's inequality applied
to the reduced truncation, whose denominator is at most \(Q_N\), would
give a lower bound \(c_L/Q_N^d\). This again contradicts the upper
bound as \(N\to\infty\).
\end{proof}

See \cite[Chapter XI, \S11.7]{hardy-wright}.
\chapter{Pell Equations}
\label{ch:nt-pell}

Fix a positive nonsquare integer \(D\) and let \(a_0=\lfloor\sqrt
D\rfloor\). In particular \(D\ge2\). We use the complete quotients and
convergent indices of Chapter~\ref{ch:nt-cf}. The arithmetic here takes
place in \(\mathbb Z[\sqrt D]\); no assumption of squarefreeness, and no
identification with the full ring of integers of the quadratic field,
is needed.

\section{Return points and the two norms}

\begin{lemma}[Return points of the square-root expansion]
\label{lem:nt-pell-return}
There is a least \(\ell\ge1\) such that
\[
 \sqrt D=[a_0;\overline{a_1,\ldots,a_{\ell-1},2a_0}].
\]
With \(\alpha_j=(P_j+\sqrt D)/Q_j\) the complete quotients,
\(Q_0=1\), and \(Q_j=1\) for \(j\ge1\) exactly when \(\ell\mid j\).
The convergents \(p_n/q_n\) satisfy
\begin{equation}\label{eq:nt-pell-norm}
 p_n^2-Dq_n^2=(-1)^{n+1}Q_{n+1}\qquad(n\ge0).
\end{equation}
\end{lemma}
\begin{proof}
The first complete quotient after \(\sqrt D\) is
\(\alpha_1=(\sqrt D-a_0)^{-1}>1\), with conjugate
\((- \sqrt D-a_0)^{-1}\in(-1,0)\). Thus it is reduced and its
complete quotients are purely periodic by Theorem~\ref{nt:E08}.
They have the forms \((P_j+\sqrt D)/Q_j\) with integers \(P_j,Q_j\)
and \(Q_j>0\). In detail, \(P_0=0,Q_0=1\), and rationalizing the
reciprocal after subtraction of \(a_j\) gives
\[
 P_{j+1}=a_j Q_j-P_j,\qquad Q_{j+1}=(D-P_{j+1}^2)/Q_j.
\]
The first step gives \(P_1=a_0,Q_1=D-a_0^2>0\). Subsequently
the reduced inequalities give \(-\sqrt D<P_j<\sqrt D\) and
\(-\sqrt D<P_{j+1}<\sqrt D\): the latter follows from
\(a_jQ_j\le P_j+\sqrt D<(a_j+1)Q_j\) with strictness on the left,
and from \(a_j\ge1\) and \(P_j<\sqrt D\).
So \(Q_{j+1}>0\). It is an integer because \(Q_j\mid D-P_j^2\) and
\(P_{j+1}\equiv-P_j\bmod Q_j\); the new divisibility follows from the
displayed product relation.

The reduced number \(\sqrt D+a_0\) has digit \(2a_0\) and maps
to \(\alpha_1\). A reduced predecessor of \(\alpha_1\) is unique
by Lemma~\ref{lem:nt-reduced-step}. Thus if \(\ell\) is the period
length of \(\alpha_1\), then \(\alpha_\ell=\sqrt D+a_0\).
For \(j\ge1\), \(Q_j=1\) is equivalent to
\(\alpha_j=\sqrt D+a_0\), since reducedness then gives
\(-1<P_j-\sqrt D<0\), and the coefficients in the
\(\mathbb Q\)-basis \(1,\sqrt D\) determine \(P_j,Q_j\) uniquely.
Pure periodicity and the least return length now show exactly
when \(Q_j=1\). The last digit of each period is \(2a_0\).

To obtain the norm formula, solve the convergent tail identity
\eqref{eq:nt-cf-tail} for \(\alpha_{n+1}\):
\[
 \alpha_{n+1}
   =\frac{p_{n-1}-q_{n-1}\sqrt D}{q_n\sqrt D-p_n}.
\]
The denominator is nonzero by irrationality. On rationalizing,
the coefficient of \(\sqrt D\) is
\[
 \frac{p_{n-1}q_n-q_{n-1}p_n}{Dq_n^2-p_n^2}
 =\frac{(-1)^n}{Dq_n^2-p_n^2}.
\]
Equating it to \(1/Q_{n+1}\) gives \eqref{eq:nt-pell-norm}.
\end{proof}

\begin{theorem}[Period and Pell signs]\label{nt:D08}
Let \(\ell\) be as in Lemma~\ref{lem:nt-pell-return}. Among positive
integers \(x,y\), the solutions of \(x^2-Dy^2=\varepsilon\), where
\(\varepsilon\in\{1,-1\}\), are exactly
\[
 (x,y)=(p_{k\ell-1},q_{k\ell-1}),\qquad
 k\ge1,\quad (-1)^{k\ell}=\varepsilon .
\]
The least positive solution in the ordering by \(x+y\sqrt D\)
is therefore the one with the least such \(k\). For the plus sign
it is obtained at \(k=1\) if \(\ell\) is even and \(k=2\) if
\(\ell\) is odd. The minus sign is possible exactly when \(\ell\)
is odd, and then its least solution occurs at \(k=1\).
\end{theorem}
\begin{proof}
Each pair displayed has norm \(\varepsilon\) by
\eqref{eq:nt-pell-norm}, and it has positive coordinates because
\(p_0=a_0\ge1,\ q_0=1\), and the recurrence uses positive digits
at later steps. Conversely a solution with \(x,y>0\) is coprime,
since a common divisor would have its square divide \(1\).
Moreover \(x^2\ge D y^2-1\ge2y^2-1\ge y^2\), so
\(x/y+\sqrt D>2\). Thus
\[
 \left|\sqrt D-\frac{x}{y}\right|
 =\frac{1}{y^2(x/y+\sqrt D)}<\frac{1}{2y^2}.
\]
Legendre's criterion in Corollary~\ref{cor:nt-legendre} shows that \(x/y\) is
a convergent. The convergent is reduced with positive denominator,
so its numerator and denominator are precisely \(x,y\).
By \eqref{eq:nt-pell-norm} this forces \(Q_{n+1}=1\); the return
criterion gives \(n+1=k\ell\) and the sign condition follows.
Finally \(p_n\) strictly increases for \(n\ge0\) by the recurrence
(\(p_1=a_1a_0+1>p_0\)), and \(q_n\) does not decrease, so the
ordering by \(x+y\sqrt D\) agrees with the order of the indicated
indices.
\end{proof}

\section{The fundamental unit and all positive solutions}

For \(u=x+y\sqrt D\) with integer coefficients, write
\(\overline u=x-y\sqrt D\) and \(N(u)=u\overline u=x^2-Dy^2\).
Norm is multiplicative by conjugation. If \(N(u)=\varepsilon\) with
\(\varepsilon=\pm1\), then \(u^{-1}=\varepsilon\overline u\)
also lies in \(\mathbb Z[\sqrt D]\). Conversely, a unit has norm
an invertible integer, since its norm times the norm of its inverse
is \(1\). If such a unit satisfies \(u>1\) as a real number, then
\[
 x=\tfrac12(u+\varepsilon/u)>0,\qquad
 y=\tfrac1{2\sqrt D}(u-\varepsilon/u)>0.
\]
Conversely \(x,y>0\) gives \(u>1\).

\begin{theorem}[Fundamental positive Pell solution]\label{nt:D07}
The equation \(x^2-Dy^2=1\) has a solution with \(x,y>0\).
Let \(u_+=x_++y_+\sqrt D\) be the least unit \(>1\) of norm \(1\).
All solutions with \(x,y>0\) are obtained uniquely as
\[
 x+y\sqrt D=u_+^j,\qquad j=1,2,\ldots .
\]
Among units with positive real value and norm \(1\), all integer
powers are allowed, with exponent \(0\) giving \(1\).
\end{theorem}
\begin{proof}
Theorem~\ref{nt:D08} proves existence and identifies the least one.
Alternatively, once any solution bounds the search interval above,
each of \(x,y>0\) is bounded on that interval, so the least exists.
For \(u>1\) of norm \(1\), choose the unique \(j\ge0\) such that
\(u_+^j\le u<u_+^{j+1}\). This choice exists since powers of a real
number \(>1\) strictly increase without bound. The quotient
\(w=u/u_+^j\) is an integral-coefficient unit of norm \(1\) and
satisfies \(1\le w<u_+\). If \(w>1\), it would yield a positive
integer pair contradicting minimality. Hence \(w=1\).
This gives the positive powers, and their uniqueness follows by
strict increase. Positive values below \(1\) are the inverses of
values above \(1\). The value \(1\) corresponds to the identity
since \(1,\sqrt D\) are linearly independent over \(\mathbb Q\).
\end{proof}

\begin{corollary}[All positive solutions for norm \(-1\)]
\label{cor:nt-pell-minus}
If \(\ell\) is odd and \(u_-=x_-+y_-\sqrt D>1\) is the least
unit of norm \(-1\) greater than \(1\), then
\[
 u_+ = u_-^2,\qquad
 \{x+y\sqrt D:x,y>0,\ x^2-Dy^2=-1\}
   =\{u_-u_+^j:j=0,1,2,\ldots\}.
\]
For even \(\ell\) the set of positive solutions of norm \(-1\) is
empty.
\end{corollary}
\begin{proof}
Existence and the even case follow from Theorem~\ref{nt:D08}.
When the minus solutions exist, any such \(u>1\) has \(u/u_-\ge1\)
of norm \(1\), so Theorem~\ref{nt:D07} gives the claimed
parametrization. Conversely every indicated product is \(>1\) of
norm \(-1\), hence yields positive coordinates. Also \(u_-^2\) is
\(>1\) of norm \(1\), so \(u_-^2\ge u_+\).
If the inequality were strict, \(u_+/u_-\) would have norm \(-1\)
and lie below \(u_-\). If \(u_+>u_-\), this quotient would lie above
\(1\), contradicting
the minimality of \(u_-\). If \(u_+<u_-\), then its inverse
\(u_-/u_+>1\) also has norm \(-1\) and is below \(u_-\),
again a contradiction. Equality \(u_+=u_-\) is impossible because
their norms differ. Thus \(u_+=u_-^2\).
\end{proof}

For the square-root expansion and the two Pell signs see
\cite[Sections 7.7--7.8, Theorems 7.21--7.26]{nzm}.
For example, \(D=2\) has period length \(1\), with \(u_-=1+\sqrt2\)
and \(u_+=3+2\sqrt2\). For \(D=3\) the period length is \(2\),
with \(u_+=2+\sqrt3\) and no norm-\(-1\) solution.

\backmatter
\begin{thebibliography}{10}

\bibitem{hardy-wright}
G. H. Hardy and E. M. Wright.
\emph{An Introduction to the Theory of Numbers}.
Sixth edition, revised by D. R. Heath-Brown and J. H. Silverman.
Oxford University Press, 2008.
\href{https://academic.oup.com/book/54489}{Publisher contents}.

\bibitem{apostol}
T. M. Apostol.
\emph{Introduction to Analytic Number Theory}.
Springer, 1976.
\href{https://doi.org/10.1007/978-1-4757-5579-4}{Publisher edition}.

\bibitem{nzm}
I. Niven, H. S. Zuckerman, and H. L. Montgomery.
\emph{An Introduction to the Theory of Numbers}.
Fifth edition. John Wiley \& Sons, 1991.

\bibitem{shoup}
V. Shoup.
\emph{A Computational Introduction to Number Theory and Algebra}.
Second edition. Cambridge University Press, 2009.
\href{https://www.shoup.net/ntb/}{Author's edition}.

\bibitem{lidl}
R. Lidl and H. Niederreiter.
\emph{Introduction to Finite Fields and Their Applications}.
Revised edition. Cambridge University Press, 1994.

\bibitem{gkp}
R. L. Graham, D. E. Knuth, and O. Patashnik.
\emph{Concrete Mathematics}.
Second edition. Addison-Wesley, 1994.

\bibitem{fan}
S. Fan.
\href{https://math.dartmouth.edu/~stevefan/papers/On\%20Geometric\%20Proofs\%20of\%20Theorems\%20on\%20Sums\%20of\%20Squares.pdf}
{\emph{On Geometric Proofs of Theorems on Sums of Squares}},
Sections 2 and 5.

\bibitem{ankeny}
N. C. Ankeny.
Sums of three squares.
\emph{Proceedings of the American Mathematical Society}, 8 (1957), 316--319.
\href{https://doi.org/10.1090/S0002-9939-1957-0085275-8}{DOI}.

\bibitem{mathlib}
The Mathlib Community.
\emph{Mathlib}, Lean 4.35.0 snapshot,
\href{https://github.com/leanprover-community/mathlib4/tree/56a005855ea17e0b0958f73ec801ae140b0a1679}{commit \texttt{56a005855}}.
Cited for existing Lean interfaces, independently of the mathematical
proofs in this text.

\bibitem{mathgenome}
MathGenome.
\emph{Abstract Algebra} and \emph{Single-Variable Analysis}.
Library sources accompanying the MathGenome release.

\end{thebibliography}
\end{document}
